%------------------------------------------------------------------------------%
\chapter{Revisão}
\label{cap:revisao}

%------------------------------------------------------------------------------%
\section{Preliminares}

%------------------------------------------------------------------------------%
\begin{teoria}{Vocabulário}
    
    Alguns termos usados em textos matemáticos
    \begin{description}
        \item[Teorema] Uma afirmação matematicamente verdadeira
        \item[Lema] Assim como o teorema, é uma afirmação matematicamente verdadeira. Geralmente recebe esse nome por fazer parte da demonstração de um teorema.
        \item[Cololário] Também é uma afirmação matematicamente verdadeira. Recebe esse nome pois sua demonstração é ``fácil'' depois que um teorema foi demonstrado.
        \item[Axiomas] São verdades não questionadas dentro de uma teoria Matemática, são consideradas ``evidentemente'' verdadeiras.
    \end{description}

\end{teoria}

%------------------------------------------------------------------------------%
\begin{teoria}{Símbolos Matemáticos}
    Alguns símbolos usados em textos matemáticos 
    \begin{center}
        \begin{tabular}{clcl} \hline
            $=$            & \emph{Igual a}    & $2+2=4$               &                               \\
            $\Rightarrow$  & \emph{Implica em} & $x>3 \Rightarrow x>2$ & Nunca usar no lugar do igual! \\
            $\rightarrow$  & \emph{Tende para} & $x\to\infty$          & Usado no contexto de limites  \\
            $\forall$      & \emph{Para todo}  & $x^2 \geq 0 \quad \forall x\in\R$ & \\
            $\exists$      & \emph{Existe}     & $\exists  x \in\R \text{ tal que } x^2 > 4$ & \\
            $\nexists$     & \emph{Não existe} & $\quad\nexists x \in\R \text{ tal que } x^2 < 0\quad$ & \\
            $\infty$       & \emph{Infinito}   & \multicolumn{2}{c}{Infinto não é um número e não podemos fazer contas com ele} \\ \hline
        \end{tabular}
    \end{center}
    
    Os símbolos tem significados muito específicos, não troque um por outro sem ter certeza que de os significados são equivalentes. 
    Por exemplo, sob hipótese nenhuma, troque o símbolo $=$ por $\to$ ou $\Rightarrow$. 

\end{teoria}

%------------------------------------------------------------------------------%
\section{Conjuntos Numéricos e Desigualdades}

%-----------------------------------------------------------------------------%
\begin{teoria}{Conjuntos Numéricos}
    Os conjuntos numéricos mais comuns, usados em Matemática, são  \\
    \begin{tabular}{clcl}
        $\N$ & \emph{Naturais}  &               $1,\,2,\,3,\,\dots$               & Números que usamos para contar                 \\
        $\Z$ & \emph{Inteiros}  & $\dots,\text{-}2,\,\text{-}1,\,0,\,1,\,2,\dots$ & Todos os números inteiros                      \\
        $\Q$ & \emph{Racionais} &         $\dfrac{1}{3},\; \dfrac{7}{9}$          & Frações, ou razões, entre números inteiros     \\
        $\R$ & \emph{Reais}     &              $\sqrt{2},\;\pi,\;e $              & Todas as distâncias possíveis e seus negativos
    \end{tabular}
\end{teoria}    

\begin{teoria}{}
    Note que cada um desses conjuntos tem infinitos números e está contido dentro do próximo
    \[ \N \subset \Z \subset \Q \subset \R \]
\end{teoria}    

\begin{teoria}{}
    Um número real que não é um número racional é chamado de \textbf{Irracional}. 

    Os números irracionais não possuem representação decimal finita
    \begin{align*}
            \sqrt{2} & =\,1,41421356237309514547462185873882845044136047363281\dots \\
            e        & =\,2,71828182845904509079559829842764884233474731445312\dots \\
            \pi      & =\,3,14159265358979311599796346854418516159057617187500\dots
    \end{align*}
    Nessa disciplina basta conhecer esses números com uma casa decimal, isso é,
    \[\sqrt{2}=1,4 \qquad e=2,7 \qquad \pi=3,1\]
\end{teoria}

%------------------------------------------------------------------------------%
\begin{teoria}{Operações com Conjuntos}
    
    Se $A$ e $B$ são conjuntos, podemos definir as operações 
    \begin{description}
        \item[Interseção]  Conjunto formado pelos elementos que pertencem, ao mesmo tempo, a $A$ e a $B$
        \[A\cap B\]
        \item[União]       Conjunto formado pelos elementos que pertencem a $A$ ou a $B$
        \[A\cup B\]
        \item[Complemento] Conjunto formado pelos elementos de $A$ que não pertencem a $B$
        \[A\setminus B\]
    \end{description}
\end{teoria}

%------------------------------------------------------------------------------%
\begin{teoria}{Reta Real}
    
    Podemos pensar nos números reais, $\R$ como sendo todas as distâncias possíveis e seus negativos. 
    Representamos esses números em uma reta, que chamamos de \textbf{Reta Real}.
    
    \addgraphic{0.5}{reta-real}
    
\end{teoria}

%------------------------------------------------------------------------------%
\begin{teoria}{Intervalos}
    
    Em geral vamos trabalhar com subconjuntos dos números reais, 
    podemos indicar esses subconjuntos com a notação de conjuntos
    \[ A = \{ x\in\R \;/\; 2 < x < 4 \} \]
    \[ B = \{ x\in\R \;/\; x < 0 \text{ ou } x > 5 \} \]
    \[ C = \{ 1,\, 2,\, e,\, 3,\, \pi,\, 4 \} \]
    
    Porém, muitos desses subconjuntos são pedaços contínuos da reta real como o conjunto $A$ descrito acima. 
    Esse tipo de subconjunto dos reais é chamado de \textbf{Intervalo}.
\end{teoria}    

\begin{teoria}{}
    Os intervalos podem ser abertos ou fechados dependendo se os seus pontos extremos pertencem ou não aos intervalos.
    \begin{align*}
        [\,a,\,b\,] &\;= \{ x\in\R \st a \leq x \leq b \} \\
        (\,a,\,b\,) &\;= \{ x\in\R \st a   <  x   <  b \} \\    
        [\,a,\,b\,) &\;= \{ x\in\R \st a \leq x   <  b \} \\
        (\,a,\,b\,] &\;= \{ x\in\R \st a   <  x \leq b \}
    \end{align*}
\end{teoria}    

\begin{teoria}{}    
    Cada um desses intervalos pode ser representado graficamente marcando um segmento 
    da reta real e indicando se as extremidades pertencem ou não ao intervalo \\
    \begin{tabular}{lll}
        $[\,a,\,b\,]$  & \emph{Intervalo Fechado}     & {\includegraphics[height=4mm,width=.3\columnwidth,trim=0 30 0 40,clip]{intervalo-fechado}          } \\[5mm]
        $(\,a,\,b\,)$  & \emph{Intervalo Aberto}      & {\includegraphics[height=4mm,width=.3\columnwidth,trim=0 34 0 50,clip]{intervalo-aberto}           } \\[5mm]
        $[\,a,\,b\,)$  & \emph{Intervalo Semi-aberto} & {\includegraphics[height=4mm,width=.3\columnwidth,trim=0 30 0 40,clip]{intervalo-fechado-esquerda} } \\[5mm]
        $(\,a,\,b\,]$  & \emph{Intervalo Semi-aberto} & {\includegraphics[height=4mm,width=.3\columnwidth,trim=0 30 0 40,clip]{intervalo-fechado-direita}  } \\[5mm]
    \end{tabular}
\end{teoria}    

\begin{teoria}{}    
    Os intervalos podem ser infinitos, mas, note que, os intervalos sempre são \textbf{abertos no infinito}. 
    \[ [\,a,\,\infty\,) \qquad (\,-\infty,\,b\,] \qquad (\,-\infty,\,\infty\,)  \]
\end{teoria}    

%------------------------------------------------------------------------------%
\begin{teoria}{Desigualdades}
    Uma desigualdade são condições que envolvem relações de maior que, 
    $>$, ou menor que, $<$. Muitas vezes suas soluções são expressas em 
    termos de intervalos ou uniões de intervalos.
    \\
    
    Propriedades das Desigualdades
    \begin{propriedades}
        \item Se $a<b$ e $b<c$, então $a<c$
        \item Se $a<b$, então $a+c<b+c$
        \item Se $a<b$ e $c>0$, então $ac<bc$
        \item Se $a<b$ e $c<0$, então $ac>bc$
    \end{propriedades}
    
    Propriedades análogas são válidas se a relação $<$ entre $a$ e $b$ for trocada por $\leq$, $>$ ou $\geqslant.$
\end{teoria}    

%------------------------------------------------------------------------------%
\begin{teoria}{Valor Absoluto}
    
    O valor absoluto de um número é definido por
    \[  \abs{x} = \begin{cases}
    \phantom{-}x & \text{se } x \geq 0 \\
    -x & \text{se } x < 0
    \end{cases} \]
\end{teoria}    

\begin{teoria}{}   
    Propriedades do Valor Absoluto
    \begin{propriedades}
        \item \( \abs{-x} = \abs{x}\)
        \item \( \abs{xy} = \abs{x}\,\abs{y}\)
        \item \( \abs{\frac{x}{y}} = \frac{\abs{x}}{\abs{y}} \qquad y\neq 0\)
        \item \( \abs{x+y} \leq \abs{x}+\abs{y}\)
    \end{propriedades}
\end{teoria}    

\begin{teoria}{}   
    Se $a>0$ 
    \begin{enumerate}
        \item $\abs{x} = a\qquad$ significa o mesmo que $\qquad x=a$ ou $x=-a$
        \item $\abs{x} < a\qquad$ significa o mesmo que $\qquad -a < x < a$
        \item $\abs{x} > a\qquad$ significa o mesmo que $\qquad x > a$ ou $x<-a$
    \end{enumerate}    
\end{teoria}    

%------------------------------------------------------------------------------%
\section{Operações Aritméticas}

%------------------------------------------------------------------------------%
\begin{teoria}{Precedência de Operadores}
    Em primeiro lugar, deve-se efetuar as multiplicações e divisões, da esquerda para a direita. 
    Em seguida, são efetuadas as somas e subtrações, também da esquerda para a direita.
    Quando desejamos efetuar as operações em outra ordem, somos obrigados a usar parênteses.
\end{teoria}

%------------------------------------------------------------------------------%
\begin{teoria}{Propriedades da Soma e da Multiplicação}
    Suponha que $a$, $b$ e $c$ sejam números reais.
    \begin{propriedades}
        \item $a + b = b + a$             \tabto{6cm} Comutatividade da soma
        \item $(a + b) + c = a + (b + c)$ \tabto{6cm} Associatividade da soma
        \item $a b = b a$                 \tabto{6cm} Comutatividade da multiplicação
        \item $(ab)c = a(bc)$             \tabto{6cm} Associatividade da multiplicação
        \item $a(b + c) = ab + ac$        \tabto{6cm} Distributividade
    \end{propriedades}
\end{teoria}

%------------------------------------------------------------------------------%
\begin{teoria}{Propriedades de Números Negativos}
    Suponha que $a$ e $b$ sejam números reais.
    \begin{propriedades}
        \item $(-1)a = -a$
        \item $-(-a) = a$
        \item $(-a)b = a(-b) = -(ab)$
        \item $(-a)(-b) = ab$
        \item $-(a + b) = -a - b$
        \item $-(a - b) = -a + b = b - a$
    \end{propriedades}
\end{teoria}

%------------------------------------------------------------------------------%
\begin{teoria}{Propriedades das Frações}
    
    Assumindo que $a$, $b$, $c$ e $d$ são números reais podemos escrever as seguintes propriedades das frações.
    \\
    
    \textbf{Soma e diferença de frações com o mesmo denominador}. Se $c\neq 0$,
    \[  \frac{a}{c} + \frac{b}{c} = \frac{a+b}{c}
        \qquad
        \frac{a}{c} - \frac{b}{c} = \frac{a-b}{c}
    \]
    
    \textbf{Produto de frações}. Se $b\neq 0$ e $d\neq 0$,
    \[ \frac{a}{b} \, \frac{c}{d} = \frac{a\,c}{b\,d}
    \]
    
    \textbf{Divisão de frações}. Se $b\neq 0$, $c\neq 0$ e $d\neq 0$, 
    \[ \dfrac{\,\dfrac{\,a\,}{b}\,}{\dfrac{\,c\,}{d}} = \frac{a}{b}\,\frac{d}{c} = \frac{a\,d}{b\,c}\]
    
    \textbf{Soma e diferença de frações com denominadores diferentes}. Se $b\neq 0$ e $d\neq 0$,
    \[  \frac{a}{b} + \frac{c}{d} = \frac{ad+cb}{bd}
        \qquad
        \frac{a}{b} - \frac{c}{d} = \frac{ad-cb}{bd}
    \]
     
    \textbf{Simplificação de uma fração}. Se $b\neq 0$ e $d\neq 0$,
    \[ \frac{a\,d}{b\,d} = \frac{a}{b} \]
    
\end{teoria}

%------------------------------------------------------------------------------%
\section{Expoentes e Radicais}

%------------------------------------------------------------------------------%
\begin{teoria}{Expoentes e Radicais}
    
    Se $b$ é um número real qualquer e $n$ um número inteiro positivo, a expressão $b^n$ representa
    \[  b^n = \underbrace{\;b\;b\;b\;\cdots\;b\;}_{n \text{ fatores}} \]
    $b$ é chamado base e $n$ é chamado potência. 
    \\
    
    Na definição original $n$ é restrito a ser um número inteiro positivo, mas podemos estender essa definição para todos os reais.
    Para $n=0$ impomos a condição $b\neq 0$ e definimos 
    \[b^0=1\]
    Para valores negativos definimos
    \[ b^{-n} = \frac{1}{\;b^n\;} \]
    Para valores fracionários, definimos que $b^{1/n}$ como o número tal que 
    \[ (b^{1/n})^n = b \]
    se esse número existir dizemos que ele é a raiz $n$-ésima de $b$
    \[ b^{1/n} = \sqrt[n]{b\,} \]
    
\end{teoria}  

%------------------------------------------------------------------------------%
\begin{teoria}{Propriedades dos Expoentes}
    \begin{propriedades}
        \item \( a^m \times a^n = a^{m+n}\)
        \item \( \frac{a^m}{a^n} = a^{m-n}  \qquad a\neq 0\)
        \item \( (a^m)^n = a^{m\times n}\)
        \item \( (ab)^n = a^n \times b^n\)
        \item \( \left(\frac{a}{b}\right)^n = \frac{a^n}{b^n} \qquad b\neq 0\)
    \end{propriedades}
\end{teoria}  

%------------------------------------------------------------------------------%
\begin{teoria}{Propriedades na Forma de Raízes}
    Algumas dessas expressões podem ser escritas na forma de raízes
    \begin{propriedades}
        \item \( \sqrt[n]{ab} = \sqrt[n]{a} \sqrt[n]{b} \)
        \item \( \sqrt[n]{\frac{\;a\;}{b}\,} = \frac{\sqrt[n]{a\,}}{\sqrt[n]{b\,}} \)
        \item \( a^{m/n} = \sqrt[n]{a^m\,}= \big(\sqrt[n]{a\,}\,\big)^m \)
        \item \( \sqrt[m]{\sqrt[n]{a\,}\,} = \sqrt[mn]{a\,} \)
        \item \( \sqrt[n]{a^n\,} = \begin{cases}
                    \;a,     & \text{se } n \text{ é ímpar} \\
                    \abs{a}, & \text{se } n \text{ é par}
                 \end{cases} \)
    \end{propriedades}
\end{teoria}

%------------------------------------------------------------------------------%
\begin{teoria}{Logaritmos}
    Seja $a$ uma constante real tal que $a > 0$ e $a\neq 1$. Se $x > 0$, então dizemos que
    \[ y = \log_a(x) \qquad \text{ se e somente se } \qquad a^y = x \]
    Dizemos que $y = \log_a(x)$ é o logaritmo de $x$ na base $a$.
\end{teoria}

%------------------------------------------------------------------------------%
\begin{teoria}{Propriedades do Logaritmo}
    
    Propriedades derivadas da definição de logaritmo
    \[
    \begin{array}{ccc}
    	\textbf{Propriedade} &                        \textbf{Motivo}                        &    \textbf{Exemplo}     \\[2mm]
    	   \log_a(1) = 0     &                  \text{Sabemos que } a^0 = 1                  &      \log_8(1) = 0      \\[2mm]
    	   \log_a(a) = 1     &                  \text{Sabemos que } a^1 = a                  &      \log_3(3) = 1      \\[2mm]
    	  \log_a(a^x) = x    &             \log_a(x) \text{ é a inversa de } a^x             &    \log_7(7 4 ) = 4     \\[2mm]
    	 a^{\log_a(x)} = x   & \null\qquad a^x \text{ é a inversa de } \log_a(x) \qquad\null & 10^{\log_{10}(13)} = 13
    \end{array}
    \]
    
    Seja $a$ uma constante real tal que $a > 0$ e $a\neq 1$, e seja $c$ uma constante real qualquer. \\
    Se $x > 0$ e $y > 0$, então,
    \begin{propriedades}
        \item Logaritmo do produto   \tabto{5cm}  $\displaystyle \log_a(xy) = \log_a(x) + \log_a(y) $
        \item Logaritmo do quociente \tabto{5cm}  $\displaystyle \log_a\left(\frac{x}{y}\right) = \log_a(x) - \log_a(y) $
        \item Logaritmo da potência  \tabto{5cm}  $\displaystyle \log_a(x^c) = c\log_a(x) $
    \end{propriedades}
\end{teoria}

%------------------------------------------------------------------------------%
\begin{teoria}{Logaritmos Usuais}
    Os logaritmos mais comumente empregados possuem uma notação particular, para facilitar seu uso. São eles:
    \\
    
    O logaritmo na base 10, também chamado \textbf{Logaritmo Comum ou Decimal}, que é apresentado sem a indicação da base.
    \[\log(x) = \log_{10}(x) \]
        
    O logaritmo na base $e\approx 2,71$, também chamado \textbf{Logaritmo Natural ou Neperiano}, que é representado por $\ln$.
    \[\ln(x) = \log_e(x)\]

\end{teoria}

%------------------------------------------------------------------------------%
\section{Expressões Algébricas}

%------------------------------------------------------------------------------%
\begin{teoria}{Expressões Algébricas}
    
    Expressões Algébricas são expressões que envolvem apenas as quatro operações básicas, expoentes e radicais.
    
    Propriedade Distributiva dos Números Reais
    \begin{propriedades}
        \item \( ab + ac = a\;(b+c) \)
    \end{propriedades}
    
    Algumas identidades úteis para manipulações algébricas
    \begin{enumerate}
        \item \( (a+b)^2 = a^2 + 2ab + b^2 \)
        \item \( (a-b)^2 = a^2 - 2ab + b^2 \)
        \item \( a^2 - b^2 = (a+b)(a-b) \)
        \item \( a^3 + b^3 = (a+b)(a^2-ab+b^2) \)
        \item \( a^3 - b^3 = (a-b)(a^2+ab+b^2) \)
    \end{enumerate}
\end{teoria}

%------------------------------------------------------------------------------%
\begin{teoria}{Fatoração}
    \textbf{Fatoração} é o processo de decompor uma expressão algébrica como o produto de outras expressões algébricas.
\end{teoria}
    
%------------------------------------------------------------------------------%
\begin{exemplos}
    \begin{exemplo}
        Um exemplo do uso da propriedade distributiva para fatorar uma expressão, é a relação
    \[ 3x^2 - x = x(3x-1)  \]    
    \end{exemplo}
\end{exemplos}

%------------------------------------------------------------------------------%
\begin{teoria}{Polinômios}
    Se $n$ é um número inteiro positivo, o termo $x^n$ é chamado monômio. 
    Um polinômio é uma expressão da forma
    \[ a_nx^n + a_{n-1}x^{n-1} + a_{n-2}x^{n-2} + \cdots + a_1x + a_0 \]
    onde $n$ é um inteiro não negativo e é chamado de grau do polinômio, 
    os coeficientes $a_i$ são números reais, com~$a_n\neq 0$.
    \\
    
    Raiz de um polinômio é um número $x$ que faça com que o valor do polinômio seja zero
    \[ a_nx^n + a_{n-1}x^{n-1} + a_{n-2}x^{n-2} + \cdots + a_1x + a_0 = 0 \]
\end{teoria}    

%------------------------------------------------------------------------------%
\begin{teoria}{Polinômio do Segundo Grau}
    Um polinômio do segundo grau tem a seguinte forma
    \[ ax^2 + bx + c \]
    onde $a\neq 0$. 
\end{teoria}

%------------------------------------------------------------------------------%
\begin{teoria}{Raízes de um Polinômio do Segundo Grau}
    Suas raízes podem ser calculadas pela fórmula de Bhaskara
    \[ x = \frac{-b\pm\sqrt{\Delta\,}}{2a} \qquad \Delta = b^2 -4ac \]
    Se $\Delta < 0$ o polinômio não tem raízes reais, se $\Delta=0$ ele tem duas raízes iguais.
    \\
    
    Se o polinômio de segundo grau possui duas raízes reais distintas, ele pode ser fatorado como 
    \[ p = ax^2 + bx + c = a(x-x_1)(x-x_2)\]
    Se ele possui duas raízes reais iguais, ele é fatorado como 
    \[ p = ax^2 + bx + c = a(x-x_1)^2\]
\end{teoria}

%------------------------------------------------------------------------------%
\begin{teoria}{Coordenadas do Vértice da Parábola}
    O vértice da parábola pode ser calculado pelas fórmulas
    \[x_v = \frac{x_1+x_2}{2} = -\frac{b}{2a} \qquad y_v = f(x_v) = -\frac{\Delta}{4a}\]
\end{teoria}

%------------------------------------------------------------------------------%
\begin{teoria}{Inequações Quadráticas}
    
    Para determinar quais valores de $x$ real satisfazem cada uma dessas inequações
    \[ p = ax^2 + bx + c > 0  \qquad \text{ ou }\qquad  p = ax^2 + bx + c < 0  \]
    \begin{enumerate}
        \item Calculamos as raízes do polinômio do segundo grau $x_1$ e $x_2$
        \begin{enumerate}
            \item Se não existirem raízes reais, o polinômio é sempre maior ou sempre menor do que zero
            \item Se as raízes foram reais e iguais, o polinômio nunca muda de sinal, apenas se anula na raiz
        \end{enumerate}
        \item Se o polinômio tem duas raízes reais distintas, fatoramos o polinômio como 
            \[ p = ax^2 + bx + c = a(x-x_1)(x-x_2)\]
        \item Construímos a tabela 
            \[\begin{array}{|c|c|c|c|} \hline
            	\multicolumn{4}{|l|}{\hspace{28mm} x_1 \hspace{12mm} x_2}      \\ \hline
                (x-x_1) & + \text{ ou } -  & + \text{ ou } - & + \text{ ou } - \\
                (x-x_2) & + \text{ ou } -  & + \text{ ou } - & + \text{ ou } - \\ \hline
                   p    & + \text{ ou } -  & + \text{ ou } - & + \text{ ou } - \\ \hline
            \end{array} 
            \]
        \item Podemos agora determinar para quais valores de $x$ o polinômio $p$ é positivo ou negativo.
    \end{enumerate}
    
\end{teoria}

%------------------------------------------------------------------------------%
\section{Geometria}

%------------------------------------------------------------------------------%
\begin{teoria}{Triângulos e Círculos}
    
    Em um triângulo retângulo vale o Teorema de Pitágoras
    \[ a^2 = b^2 + c^2 \]
    onde $a$ é a medida da hipotenusa e $b$ e $c$ são as medidas dos catetos.

    A área de um triângulo de base $b$ e altura $h$ é \[A = \frac{bh}{2}\]
        
    A área de um círculo de raio $r$ é \[A=\pi r^2\]
    O perímetro de um círculo de raio $r$ é \[P=2\pi r\]
        
\end{teoria}    

%------------------------------------------------------------------------------%
\begin{teoria}{Sólidos}
    
    O volume, $V$, e a área, $A$, de uma esfera de raio $r$ são
    \[ V = \frac{4}{3} \pi r^3    \;\;\qquad A = 4\pi r^2 \]

    O volume, $V$, e a área, $A$, de um cilindro de raio $r$ e altura $h$ são
    \[ V = \pi r^2 h              \;\;\qquad A = 2\pi r^2 + 2\pi rh\]
    
    O volume, $V$, e a área, $A$, de um cone de raio $r$ e altura $h$ são
    \[ V = \frac{1}{3}\pi r^2 h   \qquad   A = \pi r \sqrt{r^2+h^2} \]
    
\end{teoria}    

%------------------------------------------------------------------------------%
%------------------------------------------------------------------------------%
\section{Exercícios}

\openAnswers{02-Revisao/02-Revisao-Exercicios-ans}

\vspace{\baselineskip}
%------------------------------------------------------------------------------%
% \chapter{Revisão}
%------------------------------------------------------------------------------%

%------------------------------------------------------------------------------%
\begin{exercicios}{Conjuntos Numéricos e Desigualdades}
    
    \begin{exercicio}
        Determine se as afirmações são verdadeiras ou falsas.
        \begin{mathlist}[2]
            \item -3 < -20
            \item  \frac{2}{3} >  \frac{5}{6}
            \item -\frac{5}{6} < -\frac{11}{12}
        \end{mathlist}
        \begin{answers}
            \begin{alphalist}[3]
            \item Falso
            \item Falso
            \item Falso
        \end{alphalist}
        \end{answers}
    \end{exercicio}
        
    \begin{exercicio}
        Represente os intervalos dados na reta real.
        \begin{mathlist}[2]
            \item (3, 6)
            \item [-1, 4)
            \item \left[ - \frac{6}{5}, - \frac{1}{2} \right] 
            \item (- \infty, 5]
        \end{mathlist}
    \end{exercicio}
        
    \begin{exercicio}
        Encontre os valores de $x$ que satisfazem as desigualdades.
        \begin{mathlist}[2]
            \item 2x + 4 < 8
            \item -6 > 4 + 5x
            \item -4x \geq 20
            \item -6 < x - 2 < 4
            \item 0 \leq x + 1 \leq 4
            \item x + 1 > 4 \) ou $x + 2 < -1$
            \item x - 4 \leq 1 \) e $x + 3 > 2$
        \end{mathlist}
    \end{exercicio}
    
    \begin{exercicio}
        Descreva os intervalos abaixo usando desigualdades e desenhe-os na reta real.
        \begin{mathlist}[2]
            \item (-2,0)
            \item [1,6)
            \item (-3,\infty)
            \item (-\infty; 12,5]
            \item [-4,5]
            \item (-5, -2]
        \end{mathlist}    
    \end{exercicio}
    
    \begin{exercicio}
        Escreva os conjuntos abaixo na forma de intervalos e desenhe-os na reta real.
        \begin{mathlist}
            \item \{ x \in \R \st x \geq 0,17 \}
            \item \{ x \in \R \st x \geq 4 \}
            \item \{ x \in \R \st - 3 < x < -1 \}
            \item \{ x \in \R \st - 1 \leq x \leq 0 \}
            \item \{ x \in \R \st 1/100 \leq x < 100 \} 
            \item \{ x \in \R \st x \leq -2 \text{ ou } x > 5 \}
            \item \{ x \in \R \st x > -4 \}
            \item \{ x \in \R \st x < 3 \}
        \end{mathlist}
    \end{exercicio}
    
    \begin{exercicio}
        Considerando os conjuntos
        \[ A = \{x \in \R \st x \geq 1 \} \]
        \[ B = \{x \in \R \st x < 2 \} \]
        \[ C = \{x \in \R \st - 2 < x \leq 4 \} \]
        determine
        \begin{mathlist}[2]
            \item A\cup C
            \item B\cup C
            \item A\cup B
            \item A\cap C
            \item B\cap C
            \item A\cap B
        \end{mathlist}
    \end{exercicio}

    \begin{exercicio}
        Reescreva os intervalos abaixo na forma mais simples e compacta possível.
        \begin{mathlist}[2]
            \item [-4,5] \cup [-2,1]
            \item [-2,3] \cup [3,4]
            \item (0,2) \cup [2,8]
            \item (-1,4) \cup (4,6)
            \item [-4,1) \cup (-2,2]
            \item [-4, -52) \cup (-3, -2)
            \item (-\infty,6) \cup (4,7)
            \item (-\infty,5) \cup (-1,\infty)
        \end{mathlist}
    \end{exercicio}
    
    \begin{exercicio}
        Escreva os conjuntos abaixo usando desigualdades e represente-os na reta real.
        \begin{mathlist}[2]
            \item (-3,1) \cup (-1,2)
            \item [-2,2) \cap ( 1/2,4]
            \item [ 1,4) \cup (1,6]
            \item (-\infty, 2] \cap (-2,0]
            \item (-\infty,-2] \cup [3,\infty)
            \item (-\infty, 8) \cap (8,\infty)
        \end{mathlist}
    \end{exercicio}
    
    \begin{exercicio}
        Descreva os conjuntos abaixo usando a notação de intervalo.
        \addgraphic{.8}{exercicio-intervalos-01}
    \end{exercicio}
    
    \begin{exercicio}
        Considerando os conjuntos
        \[
        A = (-\infty, -3] \hfill
        B = (-1,7)        \hfill
        C = [-5, 6]       \hfill\null 
        \]
        determine
        \begin{mathlist}[2]
            \item  A \cup C
            \item  B \cup C
            \item  A \cap C
            \item  B \cap C
            \item  A \cup B \cup C
            \item  A \cap B \cap C
            \item  (A \cup B) \cap C
            \item  A \cup (B \cap C)
            \item  (A \cap B) \cup C
        \end{mathlist}
    \end{exercicio}
        
    

% 1
\begin{exercicio}
    Resolva as inequações.
    \begin{mathlist}[2]
        \item 2x > 3
        \item 8x \geq -5
        \item x - 4 \leq 5
        \item \frac{a}{2} < 7
        \item 3z - \frac{1}{2} > \frac{1}{4}
        \item x + 1 \geq -1
        \item -x \leq 6
        \item 3 \geq -9x
        \item -\frac{w}{4} > 8
        \item 3 -2y < 7
    \end{mathlist}
    \begin{answers}
        \begin{mathlist}[3]
            \item x \in \left(\sfrac{3}{2}, \infty\right)
            \item x \in \left[- \sfrac{5}{8}, \infty\right)
            \item x \in \left(-\infty, 9\right]
            \item a \in \left(-\infty, 14\right)
            \item z \in \left(\sfrac{1}{4}, \infty\right)
            \item x \in \left[-2, \infty\right)
            \item x \in \left[-6, \infty\right)
            \item x \in \left[- \sfrac{1}{3}, \infty\right)
            \item w \in \left(-\infty, -32\right)
            \item y \in \left(-2, \infty\right)
        \end{mathlist}
    \end{answers}
\end{exercicio}

% 2
\begin{exercicio}
  Resolva as inequações.
  \begin{mathlist}[1]
    \item 1 - 2(x - 1) < 2
    \item 2 - 3x \geq x + 14
    \item 5v - 32 \leq 4 - 7v
    \item 2 - z > 3(z + 3)
    \item 2(3x + 1) < 4(5 - 2x)
    \item 8(x + 3) > 12(1 - x)
    \item \frac{2}{3} - \frac{1}{2} x \geq \frac{1}{6} + x
    \item 3(3x - 2) + 2 \left(x + \frac{1}{2} \right) \leq 19 - x
    \item \frac{3x}{2} + \frac{x}{3} + \frac{x}{6} > 0
    \item \frac{1}{3} + \frac{x}{2} < \frac{5}{6} - \frac{2x}{3}
    \item \frac{3x+1}{4} - 1 \geq \frac{1}{2} - 2x
    \item \frac{1-2x}{3} + \frac{x-2}{6} > \frac{x+3}{2} - 1
    \item \frac{2}{5}x + 1 \leq \frac{1}{5} - 2x
    \item \frac{x+2}{3} + \frac{2-3x}{2} < \frac{4x}{3}
    \item \frac{x}{3} - \frac{x+1}{2} < \frac{1-x}{4}
    \item 3(1 - 2x) < 2(x + 1) + x - 7
    \item \frac{x+10}{5} > -x + 6
    \item \frac{3x-1}{4} + \frac{1-4x}{2} < 1
    \item \frac{2x}{5} + 1 \leq 2 \left(x + \frac{3}{5} \right)
    \item \frac{1-2x}{3} - \frac{1+3x}{2} \geq 2
  \end{mathlist}
  \begin{answers}
      \begin{mathlist}[3]
        \item x \in \left(\sfrac{1}{2}, \infty\right)
        \item x \in \left(-\infty, -3\right]
        \item v \in \left(-\infty, 3\right]
        \item z \in \left(-\infty, - \sfrac{7}{4}\right)
        \item x \in \left(-\infty, \sfrac{9}{7}\right)
        \item x \in \left(- \sfrac{3}{5}, \infty\right)
        \item x \in \left(-\infty, \sfrac{1}{3}\right]
        \item x \in \left(-\infty, 2\right]
        \item x \in \left(0, \infty\right)
        \item x \in \left(-\infty, \sfrac{3}{7}\right)
        \item x \in \left[\sfrac{5}{11}, \infty\right)
        \item x \in \left(-\infty, - \sfrac{1}{2}\right)
        \item x \in \left(-\infty, - \sfrac{1}{3}\right]
        \item x \in \left(\sfrac{2}{3}, \infty\right)
        \item x \in \left(-\infty, 9\right)
        \item x \in \left(\sfrac{8}{9}, \infty\right)
        \item x \in \left(\sfrac{10}{3}, \infty\right)
        \item x \in \left(- \sfrac{3}{5}, \infty\right)
        \item x \in \left[- \sfrac{1}{8}, \infty\right)
        \item x \in \left(-\infty, -1\right]
    \end{mathlist}
  \end{answers}
\end{exercicio}

% 3
\begin{exercicio}
    Resolva as inequações e represente suas soluções na reta real.
    \begin{mathlist}[2]
        \item 1 < 2x < 3
        \item -3 \leq 4x \leq 8
        \item -1 \leq x + 2 \leq 5
        \item 0 \leq 2x - 2 \leq 6
        \item -6 \leq -2(x - 1) \leq 0
        \item 2 \leq \frac{x}{3} < 4
        \item -3 < \frac{3x}{2} \leq 6
        \item -\frac{1}{4} \leq \frac{3x-4}{7} \leq \frac{1}{2}
        \item \frac{1}{6} < \frac{2x-13}{12} < \frac{2}{3}
        \item -6 \leq 15-3(4x+7) \leq 18
        \item -4 \leq \frac{5x-4}{6} \leq 1
        \item -1 \leq \frac{4s-3}{5} \leq \frac{3}{15}
        \item -9 \leq \frac{5-8x}{3} \leq 1
        \item -\frac{5}{4} \leq \frac{2x-3}{2} \leq \frac{7}{2}
        \item -\frac{3}{2} \leq \frac{2x-5}{3} \leq \frac{1}{6}
    \end{mathlist}
    \begin{answers}
        \begin{mathlist}[3]
            \item x \in \left(\sfrac{1}{2}, \sfrac{3}{2}\right)
            \item x \in \left[- \sfrac{3}{4}, 2\right]
            \item x \in \left[-3, 3\right]
            \item x \in \left[1, 4\right]
            \item x \in \left[1, 4\right]
            \item x \in \left[6, 12\right)
            \item x \in \left(-2, 4\right]
            \item x \in \left[\sfrac{3}{4}, \sfrac{5}{2}\right]
            \item x \in \left(\sfrac{15}{2}, \sfrac{21}{2}\right)
            \item x \in \left[-2, 0\right]
            \item x \in \left[-4, 2\right]
            \item s \in \left[- \sfrac{1}{2}, 1\right]
            \item x \in \left[\sfrac{1}{4}, 4\right]
            \item x \in \left[\sfrac{1}{4}, 5\right]
            \item x \in \left[\sfrac{1}{4}, \sfrac{11}{4}\right]
        \end{mathlist}
    \end{answers}
\end{exercicio}

% 4
\begin{exercicio}
    Resolva as desigualdades e represente a solução na reta real.
    \begin{mathlist}[2]
        \item \left(x - 2\right) \left(x - 4\right) \geq 0
        \item \left(x + 1\right) \left(x - 3\right) \leq 0
        \item x \left(2 x - 1\right) \geq 0
        \item 2 x \left(x - \frac{1}{4}\right) \leq 0
        \item -3 \left(x + 2\right) \left(x - 3\right) < 0
        \item \left(x + 3\right) \left(3-5x\right) \geq 0
        \item \left(x - \frac{1}{2}\right) \left(2 x + 5\right) \leq 0
        \item \left(x - 6\right)^{2} > 0
    \end{mathlist}
    \begin{answers}
        \begin{mathlist}[2]
            \item x \in \left(-\infty, 2\right] \cup \left[4, \infty\right)
            \item x \in \left[-1, 3\right]
            \item x \in \left(-\infty, 0\right] \cup \left[\sfrac{1}{2}, \infty\right)
            \item x \in \left[0, \sfrac{1}{4}\right]
            \item x \in \left(-\infty, -2\right) \cup \left(3, \infty\right)
            \item x \in \left[-3, \sfrac{3}{5}\right]
            \item x \in \left[- \sfrac{5}{2}, \sfrac{1}{2}\right]
            \item x \in \left(-\infty, 6\right) \cup \left(6, \infty\right)
        \end{mathlist}
    \end{answers}
\end{exercicio}

% 5
\begin{exercicio}
    Resolva as desigualdades e represente a solução na reta real.
    \begin{mathlist}[2]
        \item x^2 - 3x \geq 0
        \item 3x^2 \leq 5x
        \item x^2 - 8 \leq 0
        \item x^2 + 6x \leq 0
        \item 3x^2 - \sqrt{5\,}x \geq 0
        \item x^2 + 2x > 3
        \item 49x^2 \leq 9
        \item -x^2 + 5 \leq 0
        \item -2x^2 + x \geq -6
        \item x^2 + 4x + 7 \leq 0
        \item x^2 + 2x + 1 \leq 0
        \item 2x^2 \geq 20 - 6x
        \item x^2 + 9x + 18 \leq 0
        \item x^2 - 6x + 9 \geq 0
        \item -3x^2 + 16x \leq 5
        \item 16x^2 + 25 \leq 0
        \item -4x^2 + 12x - 9 \leq 0
        \item 3x^2 \leq 2x + 5
        \item -2x^2 + 8x + 24 \leq 0
        \item -x^2 + 20x - 36 \geq 0
        \item 2x^2 - 5x \geq 3
        \item (x - 6)^2 \geq 4
    \end{mathlist}
    \begin{answers}
        \begin{mathlist}[2]
            \item x \in \left(-\infty, 0\right] \cup \left[3, \infty\right)
            \item x \in \left[0, \sfrac{5}{3}\right]
            \item x \in \left[- 2 \sqrt{2}, 2 \sqrt{2}\right]
            \item x \in \left[-6, 0\right]
            \item x \in \left(-\infty, 0\right] \cup \left[\sfrac{\sqrt{5}}{3}, \infty\right)
            \item x \in \left(-\infty, -3\right) \cup \left(1, \infty\right)
            \item x \in \left[- \sfrac{3}{7}, \sfrac{3}{7}\right]
            \item x \in \left(-\infty, - \sqrt{5}\right] \cup \left[\sqrt{5}, \infty\right)
            \item x \in \left[- \sfrac{3}{2}, 2\right]
            \item x \in \emptyset
            \item x \in \left\{-1\right\}
            \item x \in \left(-\infty, -5\right] \cup \left[2, \infty\right)
            \item x \in \left[-6, -3\right]
            \item x \in \left(-\infty, \infty\right)
            \item x \in \left(-\infty, \sfrac{1}{3}\right] \cup \left[5, \infty\right)
            \item x \in \emptyset
            \item x \in \left(-\infty, \infty\right)
            \item x \in \left[-1, \sfrac{5}{3}\right]
            \item x \in \left(-\infty, -2\right] \cup \left[6, \infty\right)
            \item x \in \left[2, 18\right]
            \item x \in \left(-\infty, - \sfrac{1}{2}\right] \cup \left[3, \infty\right)
            \item x \in \left(-\infty, 4\right] \cup \left[8, \infty\right)
        \end{mathlist}
    \end{answers}
\end{exercicio}

% 6 
\begin{exercicio}
    Resolva as desigualdades e represente a solução na reta real.
    \begin{mathlist}[2]
        \item  1 \leq x^2 + 2x - 2 \leq 6
        \item -4 \leq 3x^2 - 10 \leq 2
        \item -3 \leq x^2 - 4x \leq 5
        \item -2 \leq 2x^2 + 3x + 4 \leq 3
        \item  4 \leq (x - 6)^2 \leq 9
        \item  0 \leq x^2 - x \leq 20
    \end{mathlist}
    \begin{answers}
        \begin{mathlist}[2]
            \item x \in \left[-4, -3\right] \cup \left[1, 2\right]
            \item x \in \left[-2, - \sqrt{2}\right] \cup \left[\sqrt{2}, 2\right]
            \item x \in \left[-1, 1\right] \cup \left[3, 5\right]
            \item x \in \left[-1, - \sfrac{1}{2}\right]
            \item x \in \left[3, 4\right] \cup \left[8, 9\right]
            \item x \in \left[-4, 0\right] \cup \left[1, 5\right]
        \end{mathlist}
    \end{answers}
\end{exercicio}

% 7 
\begin{exercicio}
    Resolva as inequações e represente a solução na reta real.
    \begin{mathlist}[2]
        \item \frac{x-2}{x+3}    \leq 0
        \item \frac{x+4}{x-2}    \geq 0
        \item \frac{2x-3}{x-1}   \leq 0
        \item \frac{4x+5}{x+2}   \geq 0
        \item \frac{x-3}{2x+6}   \leq 0
        \item \frac{4-5x}{2x-1}  \geq 0
        \item 3 - \frac{x}{x+2}  \leq 0
        \item \frac{3x-2}{5-2x}  \geq 0
        \item \frac{5x}{x-4}     \geq 10
        \item \frac{2x-7}{x-2}   \geq 3
        \item \frac{3x+10}{2x-5} \geq -3
        \item \frac{3x-4}{1-6x}  \leq 2
        \item \frac{6-x}{x-4}    \geq 1
    \end{mathlist}
    \begin{answers}
        \begin{mathlist}[2]
            \item x \in \left(-3, 2\right]
            \item x \in \left(-\infty, -4\right] \cup \left(2, \infty\right)
            \item x \in \left(1, \sfrac{3}{2}\right]
            \item x \in \left(-\infty, -2\right) \cup \left[- \sfrac{5}{4}, \infty\right)
            \item x \in \left(-3, 3\right]
            \item x \in \left(\sfrac{1}{2}, \sfrac{4}{5}\right]
            \item x \in \left[-3, -2\right)
            \item x \in \left[\sfrac{2}{3}, \sfrac{5}{2}\right)
            \item x \in \left(4, 8\right]
            \item x \in \left[-1, 2\right)
            \item x \in \left(-\infty, \sfrac{5}{9}\right] \cup \left(\sfrac{5}{2}, \infty\right)
            \item x \in \left(-\infty, \sfrac{1}{6}\right) \cup \left[\sfrac{2}{5}, \infty\right)
            \item x \in \left(4, 5\right]
        \end{mathlist}
    \end{answers}
\end{exercicio}

% 8 
\begin{exercicio}
    Resolva as inequações e represente a solução na reta real.
    \begin{mathlist}[2]
        \item \frac{x^2 -3x+2}{x^2 +x} \geq 0
        \item \frac{x+8}{x^2 +7x+12} \leq 0
        \item \frac{x^2 +5}{x^2 -4} \geq 5
        \item \frac{x^2 +6}{x^2 +1} \leq 2
        \item 1 + \frac{2}{x+1} \leq \dfrac{2}{x}
        \item \frac{4-2x}{x^2 -4} \leq 3
        \item \frac{x+6}{3x^2 +2} \geq 1
        \item \frac{x-5}{2x-5} \geq x
        \item \frac{4x-7}{x+2} \leq x - 2
        \item \frac{3x-1}{x+4} + \dfrac{x}{x-5} \leq 0
        \item \frac{x}{x+1} - \dfrac{1}{x-3} \geq 2
        \item \frac{x^2 +2x+3}{x+15} \leq 1
        \item \frac{3x^2 +2x-13}{x^2 -3x-10} \geq 2
    \end{mathlist}
    \begin{answers}
        \begin{mathlist}[2]
            \item x \in \left(-\infty, -1\right) \cup \left(0, 1\right] \cup \left[2, \infty\right)
            \item x \in \left(-\infty, -8\right] \cup \left(-4, -3\right)
            \item x \in \left[- \sfrac{5}{2}, -2\right) \cup \left(2, \sfrac{5}{2}\right]
            \item x \in \left(-\infty, -2\right] \cup \left[2, \infty\right)
            \item x \in \left[-2, -1\right) \cup \left(0, 1\right]
            \item x \in \left(-\infty, - \sfrac{8}{3}\right] \cup \left(-2, 2\right) \cup \left(2, \infty\right)
            \item x \in \left[-1, \sfrac{4}{3}\right]
            \item x \in \left(-\infty, \sfrac{5}{2}\right)
            \item x \in \left(-2, 1\right] \cup \left[3, \infty\right)
            \item x \in \left(-4, \sfrac{1}{2}\right] \cup \left[\sfrac{5}{2}, 5\right)
            \item x \in \left[- \sqrt{5}, -1\right) \cup \left[\sqrt{5}, 3\right)
            \item x \in \left(-\infty, -15\right) \cup \left[-4, 3\right]
            \item x \in \left(-\infty, -7\right] \cup \left(-2, -1\right] \cup \left(5, \infty\right)
        \end{mathlist}
    \end{answers}
\end{exercicio}

% 9 
\begin{exercicio}
    Resolva as inequações e represente a solução na reta real.
    \begin{mathlist}[2]
        \item \sqrt{x} - 8              \leq 0
        \item \sqrt{x} - 3              \geq 0
        \item \sqrt{x} + 10             \geq 0
        \item \sqrt{6 - 5x} - 4         \geq 0
        \item \sqrt{2x - 3}             \leq 5
        \item \sqrt{3x + 12} + 7        \geq 0
        \item \sqrt{x + 2}              \geq x - 4
        \item \sqrt{10 - 3x}            \leq x - 2
        \item \sqrt{8x + 9}             \leq 2x + 1
        \item \sqrt{9x^2 - 1}           \geq 2 - 3x
        \item \sqrt{8 - 4x^2} + x       \leq 3
        \item \sqrt{x^2 - 4x - 12}      \leq 3x + 2
        \item \sqrt{2x^2 + 1} - 2x + 1  \geq 0
        \item \sqrt{-x^2 + 5x + 14} - 2 \leq x
    \end{mathlist}
    \begin{answers}
        \begin{mathlist}[2]
            \item x \in \left[0, 64\right]
            \item x \in \left[9, \infty\right)
            \item x \in \left[0, \infty\right)
            \item x \in \left(-\infty, -2\right]
            \item x \in \left[\sfrac{3}{2}, 14\right]
            \item x \in \left[-4, \infty\right)
            \item x \in \left[-2, 7\right]
            \item x \in \left[3, \sfrac{10}{3}\right]
            \item x \in \left[2, \infty\right)
            \item x \in \left[\sfrac{5}{12}, \infty\right)
            \item x \in \left[- \sqrt{2}, \sfrac{1}{5}\right] \cup \left[1, \sqrt{2}\right]
            \item x \in \left[6, \infty\right)
            \item x \in \left(-\infty, 2\right]
            \item x \in \left\{-2\right\} \cup \left[\sfrac{5}{2}, 7\right]
        \end{mathlist}
    \end{answers}
\end{exercicio}


%    \begin{exercicio}
%        Resolva as inequações e represente suas soluções na reta real
%        \begin{mathlist}[2]
%            \item 1 < 2x < 3
%            \item -3 \leq 4x \leq 8
%            \item -1 \leq x + 2 \leq 5
%            \item 0 \leq 2x - 2 \leq 6
%            \item -6 \leq -2(x - 1) \leq 0
%            \item 2 \leq \frac{x}{3} < 4
%            \item -3 < \frac{3x}{2} \leq 6
%            \item -\frac{1}{4} \leq \frac{3x-4}{7} \leq \frac{1}{2}
%            \item \frac{1}{6} < \frac{2x-13}{12} < \frac{2}{3}
%            \item -6 \leq 15-3(4x+7) \leq 18
%            \item -4 \leq \frac{5x-4}{6} \leq 1
%            \item -1 \leq \frac{4s-3}{5} \leq \frac{3}{15}
%            \item -9 \leq \frac{5-8x}{3} \leq 1
%            \item -\frac{5}{4} \leq \frac{2x-3}{2} \leq \frac{7}{2}
%            \item -\frac{3}{2} \leq \frac{2x-5}{3} \leq \frac{1}{6}
%        \end{mathlist}
%        \todo{Copiar resposta do exercício 4 da pg 70 do arquivo 2 do livro do Francisco}
%    \end{exercicio}                
%                               
%    \begin{exercicio}
%        Resolva as desigualdades e represente a solução na reta real
%        \begin{mathlist}[2]
%            \item (x - 2)(x - 4) \geq 0
%            \item (x + 1)(x - 3) \leq 0
%            \item (2x - 1)x \geq 0
%            \item 2x(x - 1/4) \leq 0
%            \item -3(x + 2)(x - 3) < 0
%            \item (3 - 5x)(x + 3) \geq 0
%            \item (2x + 5) (x - 1/2 ) \leq 0
%            \item (x - 6)^2 > 0
%        \end{mathlist}
%        \todo{Copiar resposta do exercício 2 da pg 105 do arquivo 2 do livro do Francisco}
%    \end{exercicio}
%
%    \begin{exercicio}
%        Resolva as desigualdades e represente a solução na reta real
%        \begin{mathlist}[2]
%            \item x^2 - 3x \geq 0
%            \item 3x^2 \leq 5x
%            \item x^2 - 8 \leq 0
%            \item x^2 + 6x \leq 0
%            \item 3x^2 - \sqrt{5\,}x \geq 0
%            \item x^2 + 2x > 3
%            \item 49x^2 \leq 9
%            \item -x^2 + 5 \leq 0
%            \item -2x^2 + x \geq -6
%            \item x^2 + 4x + 7 \leq 0
%            \item x^2 + 2x + 1 \leq 0
%            \item 2x^2 \geq 20 - 6x
%            \item x^2 + 9x + 18 \leq 0
%            \item x^2 - 6x + 9 \geq 0
%            \item -3x^2 + 16x \leq 5
%            \item 16x^2 + 25 \leq 0
%            \item -4x^2 + 12x - 9 \leq 0
%            \item 3x^2 \leq 2x + 5
%            \item -2x^2 + 8x + 24 \leq 0
%            \item -x^2 + 20x - 36 \geq 0
%            \item 2x^2 - 5x \geq 3
%            \item (x - 6)^2 \geq 4
%        \end{mathlist}
%        \todo{Copiar resposta do exercício 4 da pg 105 do arquivo 2 do livro do Francisco}
%    \end{exercicio}
%
%    \begin{exercicio}
%        Resolva as desigualdades e represente a solução na reta real
%        \begin{mathlist}[2]
%            \item  1 \leq x^2 + 2x - 2 \leq 6
%            \item -4 \leq 3x^2 - 10 \leq 2
%            \item -3 \leq x^2 - 4x \leq 5
%            \item -2 \leq 2x^2 + 3x + 4 \leq 3
%            \item  4 \leq (x - 6)^2 \leq 9
%            \item  0 \leq x^2 - x \leq 20
%        \end{mathlist}
%        \todo{Copiar resposta do exercício 5 da pg 105 do arquivo 2 do livro do Francisco}
%    \end{exercicio}
%
%    \begin{exercicio}
%        Resolva as inequações e represente a solução na reta real
%        \begin{spacing}{2}
%            \begin{mathlist}[2]
%                \item \dfrac{x-2}{x+3}    \leq 0
%                \item \dfrac{x+4}{x-2}    \geq 0
%                \item \dfrac{2x-3}{x-1}   \leq 0
%                \item \dfrac{4x+5}{x+2}   \geq 0
%                \item \dfrac{x-3}{2x+6}   \leq 0
%                \item \dfrac{4-5x}{2x-1}  \geq 0
%                \item 3 - \dfrac{x}{x+2}  \leq 0
%                \item \dfrac{3x-2}{5-2x}  \geq 0
%                \item \dfrac{5x}{x-4}     \geq 10
%                \item \dfrac{2x-7}{x-2}   \geq 3
%                \item \dfrac{3x+10}{2x-5} \geq -3
%                \item \dfrac{3x-4}{1-6x}  \leq 2
%                \item \dfrac{6-x}{x-4}    \geq 1
%            \end{mathlist}
%        \end{spacing}
%        \todo{Copiar a resposta do exercício 1 da pg 130 do arquivo 2 do livro do Francisco}
%    \end{exercicio}
%
%    \begin{exercicio}
%        Resolva as inequações e represente a solução na reta real
%        \begin{spacing}{2}
%            \begin{mathlist}[2]
%                \item \dfrac{x^2 -3x+2}{x^2 +x} \geq 0
%                \item \dfrac{x+8}{x^2 +7x+12} \leq 0
%                \item \dfrac{x^2 +5}{x^2 -4} \geq 5
%                \item \dfrac{x^2 +6}{x^2 +1} \leq 2
%                \item 1 + \dfrac{2}{x+1} \leq \dfrac{2}{x}
%                \item \dfrac{4-2x}{x^2 -4} \leq 3
%                \item \dfrac{x+6}{3x^2 +2} \geq 1
%                \item \dfrac{x-5}{2x-5} \geq x
%                \item \dfrac{4x-7}{x+2} \leq x - 2
%                \item \dfrac{3x-1}{x+4} + \dfrac{x}{x-5} \leq 0
%                \item \dfrac{x}{x+1} - \dfrac{1}{x-3} \geq 2
%                \item \dfrac{x^2 +2x+3}{x+15} \leq 1
%                \item \dfrac{3x^2 +2x-13}{x^2 -3x-10} \geq 2
%            \end{mathlist}
%        \end{spacing}
%        \todo{Copiar a resposta do exercício 2 da pg 130 do arquivo 2 do livro do Francisco}
%    \end{exercicio}
%
%    \begin{exercicio}
%        Resolva as inequações e represente a solução na reta real
%            \begin{mathlist}
%                \item \sqrt{x} - 8                     \leq 0
%                \item \sqrt{x} - 3                     \geq 0
%                \item \sqrt{x} + 10                    \geq 0
%                \item \sqrt{6 - 5x} - 4                \geq 0
%                \item \sqrt{2x - 3}                    \leq 5
%                \item \sqrt{3x + 12} + 7               \geq 0
%                \item \sqrt{x + 2}                     \geq x - 4
%                \item \sqrt{10 - 3x}                   \leq x - 2
%                \item \sqrt{8x + 9}                    \leq 2x + 1
%                \item \sqrt{9x^2 - 1}                  \geq 2 - 3x
%                \item \sqrt{8 - 4x^2} + x              \leq 3
%                \item \sqrt{x^2 - 4x - 12}             \leq 3x + 2
%                \item \sqrt{2x^2 + 1} - 2x + 1         \geq 0
%                \item \sqrt{-x^2 + 5x + 14} - 2        \leq x
%            \end{mathlist}
%        \todo{Copiar a resposta do exercício 3 da pg 130 do arquivo 2 do livro do Francisco}
%    \end{exercicio}

    \begin{exercicio}
        Calcule $\abs{3x - 10}$ para $x = 2$ e $x = 5$.
    \end{exercicio}
    
    \begin{exercicio}
        Calcule $\abs{7 - x}$ para $x = -7$, $x = 1$, $x = 7$ e $x = 12$.
    \end{exercicio}
    
    \begin{exercicio}
        Elimine o módulo da expressão $\abs{x}/x^2$.
    \end{exercicio}
    
    \begin{exercicio} % Francisco pg 145, ex 1
        Elimine o módulo das expressões.
        \begin{mathlist}[3]
            \item \abs{8}
            \item \abs{-8}
            \item -\abs{-8}
            \item \abs{3-\pi}
            \item \abs{\pi-3}
            \item \abs{\sqrt{8}-4}
            \item \abs{-\frac{10}{5^2}}
            \item \sqrt{\abs{-4}}
        \end{mathlist}	
        \begin{answers}
            \begin{mathlist}[4]
                \item 8
                \item 8
                \item -8
                \item \pi-3
                \item \pi-3
                \item 4-\sqrt{8}
                \item \sfrac{2}{5}
                \item 2
            \end{mathlist}
        \end{answers}
    \end{exercicio}
    
    \begin{exercicio}
        Calcule os valores das expressões.
        \begin{mathlist}[2]
            \item \abs{-6 + 2}
            \item 4 + \abs{-4}
            \item \frac{\abs{-12 + 4}}{\abs{16 - 12}}
            \item \abs{\frac{0,2 - 1,4}{1,6 - 2,4}}
            \item \sqrt{3\,}\; \abs{-2} + 3\abs{-\sqrt{3\,}}
            \item \abs{\pi - 1} + 2
            \item \abs{2\sqrt{3} - 5} - \abs{\sqrt{3} - 4}
        \end{mathlist}
        \begin{answers}
            \begin{mathlist}[4]
                \item 4
                \item 8
                \item 2
                \item \sfrac{1}{4}
                \item 5\sqrt{3\,}
                \item \pi + 1
                \item 1-\sqrt{3\,}
            \end{mathlist}
        \end{answers}
    \end{exercicio}
    
    \begin{exercicio}
        Suponha que $a$ e $b$ são números reais não-nulos e que $a > b$.
        Determine se as desigualdades são verdadeiras ou falsas.
        \begin{mathlist}[2]
            \item b - a > 0
            \item \frac{a}{b} > 1
            \item a^2 > b^2
            \item -a < -b
        \end{mathlist}	
    \end{exercicio}
    
    \begin{exercicio}
        Determine se as afirmações são verdadeiras para quaisquer números reais $a$ e $b$.
        \begin{mathlist}[2]
            \item \abs{-a} = a
            \item \abs{b^2} = b^2
            \item \abs{a - 4} = \abs{4 - a}
            \item \abs{a + 1} = \abs{a} + 1
            \item \abs{a + b} = \abs{a} + \abs{b}
        \end{mathlist}    
    \end{exercicio}
        
    \begin{exercicio}
        Determine se as afirmações são verdadeiras ou falsas.
	Se verdadeiras, explique o porquê. Se falsas, justifique através de um exemplo.
        \begin{alphalist}
            \item Se $a < b$, então $a - c > b - c$
            \item $\abs{a - b} \leq \abs{b} + \abs{a}$
        \end{alphalist}        
    \end{exercicio}
                
    \begin{exercicio} % Exercício 4 da pg 145 do arquivo 2 do livro do Francisco
        Elimine o módulo das expressões escrevendo-as na forma ``por partes''
        \begin{mathlist}[3]
            \item -\abs{x}
            \item \abs{x}-5
            \item 5-\abs{x}
            \item \abs{x-5}
            \item \abs{5-x}
            \item \abs{5x+1}
            \item \abs{4-3x}
            \item \abs{x^2+7}
            \item \abs{x^2-9}
        \end{mathlist}
        \begin{answers}
            \begin{mathlist}
                \item -\abs{x} = \begin{cases}
                    -x, &\text{se } x \geq 0 \\
                    \hphantom{-}x, &\text{se } x < 0
                \end{cases}
                \item \abs{x}-5 = \begin{cases}
                    \hphantom{-}x-5, &\text{se } x \geq 0 \\
                    -x-5, &\text{se } x < 0
                \end{cases}
                \item 5-\abs{x} = \begin{cases}
                    5-x, &\text{se } x \geq 0 \\
                    5+x, &\text{se } x < 0
                \end{cases}
                \item \abs{x-5} = \begin{cases}
                    x-5, &\text{se } x \geq 5 \\
                    5-x, &\text{se } x < 5
                \end{cases}
                \item \abs{5-x} = \begin{cases}
                    x-5, &\text{se } x \geq 5 \\
                    5-x, &\text{se } x < 5
                   \end{cases}
                \item \abs{5x+1} = \begin{cases}
                	\hphantom{-}5x+1, & \text{se } x \geq \sfrac{-1}{5} \\
                	-5x-1,            & \text{se } x <    \sfrac{-1}{5}
                \end{cases}
                \item \abs{4-3x} = \begin{cases}
                    4-3x, &\text{se } x \leq \sfrac{4}{3} \\
                    3x-4, &\text{se } x > \sfrac{4}{3}
                \end{cases}
                \item \abs{x^2+7} = x^2+7
                \item \abs{x^2-9} = \begin{cases}
                    x^2-9, &\text{se } x \leq -3 \text{ ou } x \geq 3 \\
                    9-x^2, &\text{se } -3 < x 3
                \end{cases}
            \end{mathlist}
        \end{answers}
    \end{exercicio}

    \begin{exercicio} % Exercício 6 da pg 145 do arquivo 2 do livro do Francisco
        Calcule as expressões
        \begin{mathlist}[2]
            \item \abs{5\times(-3)}
            \item -3\abs{-5}
            \item \abs{\frac{-3}{-6}}
            \item \abs{\frac{5-17}{15-6}}
            \item \abs{-2}+6\abs{-5}
            \item \abs{\vphantom{\sum}-2+\abs{-5}}
        \end{mathlist}
        \begin{answers}
            \begin{mathlist}[2]
                \item \abs{5\times(-3)} = 15 
                \item -3\abs{-5} = -15 
                \item \abs{\frac{-3}{-6}} = \sfrac{1}{2}
                \item \abs{\frac{5-17}{15-6}} = \sfrac{3}{4}
                \item \abs{-2}+6\abs{-5} = 32
                \item \abs{\vphantom{\sum}-2+\abs{-5}} = 3
           \end{mathlist}
        \end{answers}        
    \end{exercicio}

    \begin{exercicio} % Exercício 8 da pg 145 do arquivo 2 do livro do Francisco
        Calcule as expressões 
        \begin{mathlist}[2]
            \item \abs{(-4x)\times(-6)}
            \item \abs{ \frac{3x}{-6}}
            \item \abs{-\frac{-3}{6x}}
            \item \abs{-4x}+\abs{8x}
            \item \abs{2x}-\abs{-2x}
            \item \abs{\frac{2x}{3}} - \frac{\abs{x}}{6}
            \item \frac{\abs{2x^2}}{\abs{-4xy}}
            \item \dfrac{\sqrt{x^2}}{\abs{x}}
            \item \sqrt{\abs{-2x^2}}
        \end{mathlist}	
        \begin{answers}
            \begin{mathlist}[2]
                \item \abs{(-4x)\times(-6)}=24\abs{x}
                \item \abs{ \frac{3x}{-6}}=\frac{\abs{x}}{2}
                \item \abs{-\frac{-3}{6x}}=\frac{1}{2\abs{x}}
                \item \abs{-4x}+\abs{8x}=12\abs{x}
                \item \abs{2x}-\abs{-2x}=0
                \item \abs{\frac{2x}{3}} - \frac{\abs{x}}{6}=\frac{\abs{x}}{2}
                \item \frac{\abs{2x^2}}{\abs{-4xy}} =\frac{\abs{x}}{2\abs{y}}
                \item \frac{\sqrt{x^2}}{\abs{x}}=1
                \item \sqrt{\abs{-2x^2}} =\abs{x}\sqrt{2}
            \end{mathlist}
        \end{answers}   
    \end{exercicio}
              
    \begin{exercicio} % exercício 7 da pg 145 do arquivo 2 do livro do Francisco
        Calcule a distância entre os pontos da reta real.
        \begin{alphalist}
            \item \(x_1 = -5\) e \(x_2 = -8\)
            \item \(x_1 = -10\) e \(x_2 = 10\)
            \item \(x_1 = 4,7\) e \(x_2 = 1,2\)
            \item \(x_1 = 2\) e \(x_2 = -12\)
        \end{alphalist}
        \begin{answers}
            \begin{mathlist}[4]
                \item 3
                \item 20
                \item 3,5
                \item 14
            \end{mathlist}
        \end{answers}   
    \end{exercicio}
        
    \begin{exercicio}
        Reescreva as frases abaixo usando equações modulares.
        \begin{alphalist}
            \item A distância entre $x$ e $ 2$ é igual a $3$.
            \item A distância entre $s$ e $-3$ é igual a $4$.
            \item A casa de minha avó e a casa de meu tio estão a \SI{5}{\kilo\meter} de distância.
            \item A farmácia e a padaria estão à mesma distância de minha casa.
        \end{alphalist}
    \end{exercicio}

    \begin{exercicio}
        Determine os pontos da reta real que estão a uma distância de $10$ unidades de $6$.
    \end{exercicio}

    \begin{exercicio}
        Determine os pontos da reta real que estão a uma distância de $2/3$ unidades de $-1$.
    \end{exercicio}

    \begin{exercicio}
        Determine o conjunto dos valores de $x$ que satisfazem a inequação $0 < x^2 - x < 20$
    \end{exercicio}

    \begin{exercicio} % exercício 12 da pg 145 do arquivo 2 do livro do Francisco
        Resolva as equações.
        \begin{mathlist}[3]
            \item \abs{x} = 4
            \item \abs{x} = -4
            \item x = \abs{-4}
            \item x = \abs{4}
            \item \abs{x} = \abs{4}
            \item \abs{x} = \abs{-4}
        \end{mathlist}
        \begin{answers}
            \begin{mathlist}[2]
                \item x=-4 \text{ ou } x=4
                \item \) Não há solução
                \item x=4
                \item x=4
                \item x=-4 \text{ ou } x=4 
                \item x=-4 \text{ ou } x=4 
            \end{mathlist}
        \end{answers}
    \end{exercicio}

    \begin{exercicio} % exercício 13 da pg 146 do arquivo 2 do livro do Francisco
        Resolva as equações.
        \begin{mathlist}[2]
            \item \abs{x-3} = 4
            \item \abs{x-\frac{1}{2}} = 2
            \item \abs{4-x} = \frac{1}{10}
            \item \abs{3x-1} =6
            \item \abs{x-2} =-1
            \item \abs{\frac{x-3}{4}} = 12
            \item \abs{5-4x} =1
            \item \abs{5x-2} =13
            \item \abs{\frac{2-3x}{4}} =5
            \item \abs{\frac{3x}{2}-1} =\frac{5}{2}
            \item \abs{5x-3} =3x+15
            \item \abs{2x-3} =5-4x
            \item \abs{6-x} +5x=7
            \item \abs{\frac{7+4x}{3}} =5+x
            \item \abs{3x+5} =\abs{x-3}
            \item \abs{2x+1} = \abs{2-5x}
            \item \abs{x-1} + \abs{x+2} =5
            \item \abs{2x-5} - \abs{x-2} = 3x+1
        \end{mathlist}
        \begin{answers}
            \begin{mathlist}[2]
                \item x=-4 \text{ ou } x=-2 \text{ ou } x=2 \text{ ou } x=4
                \item x=\sfrac{-3}{2} \text{ ou } x=\sfrac{3}{2}
                \item x=-1 \text{ ou } x=7 \text{ ou } x=3-\sqrt{2} \text{ ou } x=3+\sqrt{2}
                \item x=0 \text{ ou } x=4
                \item x=\sfrac{-2}{3} \text{ ou } x=\sfrac{-1}{2} \text{ ou } x=0 \text{ ou } x=2
                \item x=-6 \text{ ou } x=1 \text{ ou } x=2 \text{ ou } x=3
                \item x=3 \text{ ou } x=-3-\sqrt{2}
                \item x=-2 \text{ ou } x=2
                \item x=-7 \text{ ou } x=2 \text{ ou } x=3
                \item x=-\sqrt{13} \text{ ou } x=\sqrt{13}
            \end{mathlist}
        \end{answers}
    \end{exercicio}

%\todo{
%    \begin{exercicio} % exercício 14 da pg 146 do arquivo 2 do livro do Francisco
%        Resolva as equações.
%        \begin{answers}
%            \begin{mathlist}[2]
%                \item 
%            \end{mathlist}
%        \end{answers}
%    \end{exercicio}
%
%    \begin{exercicio} % exercício 15 da pg 146 do arquivo 2 do livro do Francisco
%        Resolva as equações.
%        \begin{answers}
%            \begin{mathlist}[2]
%                \item 
%            \end{mathlist}
%        \end{answers}
%    \end{exercicio}
%
%    \begin{exercicio} % exercício 17 da pg 146 do arquivo 2 do livro do Francisco
%        Resolva as desigualdades.
%        \begin{answers}
%            \begin{mathlist}[2]
%                \item 
%            \end{mathlist}
%        \end{answers}
%    \end{exercicio}
%
%    \begin{exercicio} % exercício 18 da pg 146 do arquivo 2 do livro do Francisco
%        Identifique, na reta real, os intervalos definidos pelas desigualdades.
%        \begin{answers}
%            \begin{mathlist}[2]
%                \item 
%            \end{mathlist}
%        \end{answers}
%    \end{exercicio}
%
%    \begin{exercicio} % exercício 19 da pg 146 do arquivo 2 do livro do Francisco
%       Resolva as desigualdades.
%       \begin{answers}
%           \begin{mathlist}[2]
%               \item 
%            \end{mathlist}
%        \end{answers}
%    \end{exercicio}
%}
    
\end{exercicios}

%------------------------------------------------------------------------------%

%------------------------------------------------------------------------------%
% \chapter{Revisão}
%------------------------------------------------------------------------------%

%------------------------------------------------------------------------------%
%\begin{exercicios}{Operações}
%
%    \begin{exercicio}
%    \end{exercicio}
%    
%\end{exercicios}

%------------------------------------------------------------------------------%

%------------------------------------------------------------------------------%
% \chapter{Revisão}
%------------------------------------------------------------------------------%

%------------------------------------------------------------------------------%
\begin{exercicios}{Expoentes e Radicais}
    
    \begin{exercicio}
        Calcule os valores das expressões.
        \begin{mathlist}[2]
            \item 27^{2/3}
            \item 8^{-4/3}
            \item \left(\frac{1}{\sqrt{3}}\right)^0
            \item (7^{1/2})^4
            \item \left[\left(\frac{1}{8}\right)^{1/3}\right]^{-2}
            \item \left[\left(- \frac{1}{3}\right)^2\right]^{-3}
            \item \left(\frac{7^{-5} \, 7^2}{7^{-2}}\right)^{-1}
            \item (125^{2/3})^{-1/2}
            \item \sqrt[3]{2^6}
            \item \frac{\sqrt{32}}{\sqrt{8}}
            \item \sqrt[3]{\frac{-8}{27}}
            \item \frac{16^{5/8} \, 16^{1/2}}{16^{7/8}}
            \item \frac{6^{2,5} \, 6^{-1,9}}{6^{-1,4}}
        \end{mathlist}	
        \begin{answers}
            \begin{mathlist}[3]
                \item 9
                \item \frac{1}{16}
                \item 1
                \item 49
                \item 4
                \item 729
                \item 7
                \item \frac{1}{5}
                \item 4
                \item 2
                \item \frac{2 \sqrt[3]{-1}}{3}
                \item 2
                \item 36
            \end{mathlist}
        \end{answers}
    \end{exercicio}

\todo{    
    \begin{exercicio}
        Copiar exercícios 1-6,7-12,13-18, pg 321 livro do Tan
    \end{exercicio}
}
    
    \begin{exercicio}
        Determine se as afirmações são verdadeiras ou falsas. Justifique sua decisão.	
        \begin{mathlist}[2]
            \item x^4 + 2x^4 = 3x^4
            \item 3^2 \, 2^2 = 6^2
            \item \frac{2^{4x}}{1^{3x}} = 2^{4x-3x}
            \item (2^2 \, 3^2)^2 = 6^4
            \item \frac{4^{3/2}}{2^4} = \frac{1}{2}
            \item (1,2^{1/2})^{-1/2} = 1		
        \end{mathlist}	
    \end{exercicio}
    
    \begin{exercicio}
        Reescreva as expressões utilizando apenas expoentes positivos.
        \begin{mathlist}[2]
            \item (xy)^{-2}
            \item \frac{x^{-1/3}}{x^{1/2}}
            \item \sqrt{x^{-1}} \, \sqrt{9x^{-3}}
            \item 12^0 (s + t)^{-3}
        \end{mathlist}			
    \end{exercicio}
    
    \begin{exercicio}
        Simplifique as expressões. (Assuma que $x$, $y$, $r$, $s$ e $t$ são positivos.)
        \begin{mathlist}[2]
            \item \frac{x^{7/3}}{x^{-2}}
            \item (49x^{-2})^{1/2}
            \item (x^2y^{-3})(x^{-5}y^3)
            \item \frac{5x^6 y^3}{2x^2 y^7}
            \item \frac{x^{3/4}}{x^{-1/4}}
            \item \left(\frac{x^3}{-27y^{-6}}\right)^{-2/3}
            \item \left(\frac{x^{-3}}{y^{-2}}\right)^2 \left(\frac{y}{x}\right)^4
            \item \frac{(r^n)^4}{r^{5-2n}}
            \item \sqrt[3]{x^{-2}} \, \sqrt{4x^5}
            \item \sqrt{81x^6 y^{-4}}
            \item \sqrt[3]{27r^6} \, \sqrt{s^2 t^4}
        \end{mathlist}	
    \end{exercicio}
    
    \begin{exercicio}
        Utilize o fato de que $2^{1/2} \approx 1,414$ e $3^{1/2} \approx 1,732$ para calcular
        as expressões dadas sem o auxílio de uma calculadora.	
        \begin{mathlist}[2]
            \item 2^{3/2}
            \item 8^{1/2}
        \end{mathlist}	
    \end{exercicio}
    
    \begin{exercicio}
        Utilize o fato de que $10^{1/2} \approx 3,162$ e $10^{1/3} \approx 2,154$	para calcular
        as expressões dadas sem o auxílio de uma calculadora.	
        \begin{mathlist}[3]
            \item 10^{3/2}
            \item 10^{2,5}
            \item (0,0001)^{-1/3}
        \end{mathlist}	
    \end{exercicio}
    
    \begin{exercicio}
        Racionalize o denominador das expressões.
        \begin{mathlist}[3]
            \item \frac{3}{2 \sqrt{x}}
            \item \frac{5x^2}{\sqrt{3x}}
            \item \sqrt{\frac{2x}{y}}
        \end{mathlist}	
    \end{exercicio}
    
    \begin{exercicio}
        Racionalize o denominador das expressões.
        \begin{mathlist}[2]
            \item \frac{2\sqrt{x}}{3}
            \item \frac{\sqrt[3]{x}}{24}
            \item \sqrt{\frac{2y}{x}}
            \item \frac{\sqrt[3]{x^2 z}}{y}
        \end{mathlist}
    \end{exercicio}
    
\end{exercicios}

%------------------------------------------------------------------------------%

%------------------------------------------------------------------------------%
% \chapter{Revisão}
%------------------------------------------------------------------------------%

%------------------------------------------------------------------------------%
\begin{exercicios}{Logaritmos}
    
    \begin{exercicio}
        Coloque a expressão dada na forma logarítmica
        \begin{mathlist}[2]
            \item 2^6 = 64
            \item 3^5 = 243
            \item 3^{-2} = \frac{1}{9}
            \item \left(\frac{1}{3}\right)^1 = \frac{1}{3}
            \item \left(\frac{1}{2}\right)^1{-4} = 16 
            \item 32^{3/5} = 8
            \item 81^{3/4} = 27
            \item 10^{-3} = 0,0001
        \end{mathlist}
    \end{exercicio}
    
    \begin{exercicio}
        Use o fato que $\log 3 = 0,4771$ e $\log 4 = 0,6021$ para encontrar o
        valor do logaritmo solicitado
        \begin{mathlist}[2]
            \item \log 12
            \item \log \frac{3}{4}
            \item \log 16
            \item \log \sqrt{3}
        \end{mathlist}
    \end{exercicio}
        
    \begin{exercicio}
        Utilize as propriedades dos logaritmos para simplificar as expressões dadas
        \begin{mathlist}[2]
            \item \log x(x+1)^4
            \item \log\frac{\sqrt{x+1}}{x^2+1}
            \item \ln\frac{e^x}{1+e^x}
            \item \ln xe^{-x^2}
            \item \ln x(x+1)(x+2)
            \item \ln \frac{x^{1/2}}{x^2\sqrt{1+x^2}}
            \item \ln x^x
        \end{mathlist}
        \begin{answers}
            \begin{mathlist}[2]
                \item \log x + 4\log(x+1)
                \item \frac{1}{2}\log(x+1) - \log\left(x^2+1\right)
                \item x -\ln\left(1+e^x\right)
                \item \ln x -x^2
                \item \ln x + \ln(x+1) + \ln(x+2)
                \item \frac{1}{2}\ln x - 2\ln x - \frac{1}{2}\ln\left(1+x^2\right)
                \item x\ln x
            \end{mathlist}
        \end{answers}
    \end{exercicio}
        
    \begin{exercicio}
        Use logaritmos para resolver as equações em~$t$
        \begin{mathlist}[2]
            \item e^{0,4t} = 8
            \item \frac{1}{3} e^{-3t} = 0,9
            \item 5e^{-2t} = 6
            \item 4e^{t-1} = 4
            \item \frac{50}{1+4e^{0,2t}} = 20
            \item \frac{200}{1+3e^{-0,3t}} = 100
            \item A = B e^{-t/2}
            \item \frac{A}{1+Be^{t/2}} = C
        \end{mathlist}
        \begin{answers}
          \begin{mathlist}[2]
              \item t = \frac{5}{2} \ln{\left (8 \right )}
              \item t = -\frac{1}{3}\ln(2,7)
              \item t = -\frac{1}{2}\ln{\left (\frac{6}{5} \right )}
              \item t = 1 
              \item t = 5 \ln{\left (\frac{3}{8} \right )}
              \item t = \frac{10}{3}\ln(3)
              \item t = -2\ln\left(\frac{A}{B}\right)
              \item t = 2\ln\left(\frac{A-C}{CB}\right)
            \end{mathlist}
        \end{answers}        
    \end{exercicio}
    
\end{exercicios}

%------------------------------------------------------------------------------%

%------------------------------------------------------------------------------%
% \chapter{Revisão}
%------------------------------------------------------------------------------%

%------------------------------------------------------------------------------%
\begin{exercicios}{Expressões Algébricas}
    
    % 1
\begin{exercicio}
  Simplifique as expressões abaixo, reduzindo os termos semelhantes.
  \begin{mathlist}[1]
      \item (3x + 2) + (5x - 4)
      \item (2y - 3) - (4y - 5)
      \item (y^3 - 4y^2 + y - 1) - (3y^2 + y - 6)
      \item (-5z + 2x - 6) + 3(z + 4x + 2)
      \item (2a - 5b + 3c) + (6a + 2ab - 3c)
      \item -2(a - 2b - 3ab) - 4(b + 2a - 2ab)
      \item \frac{x - 2}{2} - (2 - x)
      \item \frac{2}{3} (2x - 1) + \frac{4}{3} (2 - x)
      \item \frac{1}{2} (x + 2y - 4) + \frac{1}{3} (3y - x + 9)
      \item \frac{1}{2} (a - 3ab + 2b) + \frac{1}{3} (a - 3b + 4ab)
  \end{mathlist}
  \begin{answers}
      \begin{mathlist}[2]
        \item 8 x - 2
        \item - 2 y + 2
        \item y^{3} - 7 y^{2} + 5
        \item 14 x - 2 z
        \item 2 a b + 8 a - 5 b
        \item 2 a \left(7 b - 5\right)
        \item \frac{3 x}{2} - 3
        \item 2
        \item \frac{x}{6} + 2 y + 1
        \item \frac{a \left(- b + 5\right)}{6}
    \end{mathlist}
  \end{answers}
\end{exercicio}


% 2
\begin{exercicio}
  Expanda as expressões e simplifique-as.
  \begin{mathlist}[2]
      \item \left(\frac{x}{5}\right) \, \left(\frac{2}{3} - 2x\right)
      \item \left(- \frac{x}{2}\right) \, \left(2 - \frac{3x}{4}\right)
      \item (5x - 3) (2x + 4)
      \item (8 - 3x) (x^2 + 6)
      \item -2(1 - x) \left(3 + \frac{x}{2}\right)
      \item (0,7x - 0,2) (4 - 0,6x)
      \item \left(x - \frac{1}{2}\right) \, \left(\frac{1}{3} - x \right)
      \item \left(\frac{x}{2} - 3 \right) \, \left(\frac{5}{4} + 2x \right)
      \item \left(\frac{2x}{3} - \frac{3}{2}\right) \, \left(\frac{3}{4} - \frac{x}{3}\right)
      \item (12x - 5) (12x + 5)
      \item (3x + 4)^2
      \item (x - \sqrt{3})^2
  \end{mathlist}
  \begin{answers}
      \begin{mathlist}[2]
          \item \frac{2 x \left(- 3 x + 1\right)}{15}
          \item \frac{x \left(3 x - 8\right)}{8}
          \item 10 x^{2} + 14 x - 12
          \item - 3 x^{3} + 8 x^{2} - 18 x + 48
          \item x^{2} + 5 x - 6
          \item - 0,42 x^{2} + 2,92 x - 0,8
          \item - x^{2} + \frac{5 x}{6} - \frac{1}{6}
          \item x^{2} - \frac{43 x}{8} - \frac{15}{4}
          \item - \frac{2 x^{2}}{9} + x - \frac{9}{8}
          \item 144 x^{2} - 25
          \item 9 x^{2} + 24 x + 16
          \item x^{2} - 2 \sqrt{3} x + 3
    \end{mathlist}
  \end{answers}
\end{exercicio}


% 3
\begin{exercicio}
  Efetue os produtos abaixo.
  \begin{mathlist}[2]
      \item (x^{-1} + 3)(x + 2)
      \item (3x^2 + 2)(6 - \sqrt{x})
      \item (\sqrt{x} + 9)(\sqrt{x} - 9)
      \item \left(\frac{2}{\sqrt{x}} - 5\right) (\sqrt{x} - 1)
      \item \left(\frac{1}{x} - 1\right) \left(4 + \frac{1}{x^2}\right)
      \item 3x^2 (x^3 - 2x^2 - 4x +5)
      \item -4x^3 (x^2 + 2x - 1)
      \item xy^2 (2x + 3xy + 4y)
      \item (3x + 5) (3x^3 + 4x^2 + 2)
      \item (2 - x^2) (3x^3 + 6x^2 - x)
      \item \left(2x^2 - \frac{1}{2}\right) (x^2 + 3)
      \item (x^3 + 1) (x^4 - 3x^2 + 2)
      \item (3 - 2y + y^2) (2y^2) - 5y + 4
      \item (x - y + 1) (2x - 4y + 6)
      \item (x^2 + 2y) (3x - 2xy - y)
      \item (2x - 1)^3
      \item (x - 3) (x + 3) (x - 2)
      \item (2w - 3) (w - 1) (3w + 2)
      \item (x^2 + 3) (x^2 - 2) (2x^2 - 5)
      \item (a + 2b) (3a - b) (2a + 3b)
      \item (a - b) (a + b) (a^2 + b^2)
  \end{mathlist}
  \begin{answers}
      \begin{mathlist}[1]
          \item 3 x + 7 + \frac{2}{x}
          \item - 3 x^{\frac{5}{2}} - 2 \sqrt{x} + 18 x^{2} + 12
          \item x - 81
          \item - 5 \sqrt{x} + 7 - \frac{2}{\sqrt{x}}
          \item -4 + \frac{4}{x} - \frac{1}{x^{2}} + \frac{1}{x^{3}}
          \item 3 x^{2} \left(x^{3} - 2 x^{2} - 4 x + 5\right)
          \item 4 x^{3} \left(- x^{2} - 2 x + 1\right)
          \item x y^{2} \left(3 x y + 2 x + 4 y\right)
          \item 9 x^{4} + 27 x^{3} + 20 x^{2} + 6 x + 10
          \item x \left(- 3 x^{4} - 6 x^{3} + 7 x^{2} + 12 x - 2\right)
          \item 2 x^{4} + \frac{11 x^{2}}{2} - \frac{3}{2}
          \item x^{7} - 3 x^{5} + x^{4} + 2 x^{3} - 3 x^{2} + 2
          \item 2 y^{4} - 4 y^{3} + 6 y^{2} - 5 y + 4
          \item 2 x^{2} - 6 x y + 8 x + 4 y^{2} - 10 y + 6
          \item - 2 x^{3} y + 3 x^{3} - x^{2} y - 4 x y^{2} + 6 x y - 2 y^{2}
          \item 8 x^{3} - 12 x^{2} + 6 x - 1
          \item x^{3} - 2 x^{2} - 9 x + 18
          \item 6 w^{3} - 11 w^{2} - w + 6
          \item 2 x^{6} - 3 x^{4} - 17 x^{2} + 30
          \item 6 a^{3} + 19 a^{2} b + 11 a b^{2} - 6 b^{3}
          \item a^{4} - b^{4}
    \end{mathlist}
  \end{answers}
\end{exercicio}


    
\begin{exercicio}
  Expanda as expressões abaixo
  \begin{mathlist}[2]
            \item (x + 2)^2
            \item (3x + 8)^2
            \item (x^2 - \sqrt{5})^2
            \item (2u + 7v)^2
            \item (4 - y)^2
            \item (3 - 2y)^2
            \item (-2 - x)^2
            \item \left(\frac{x}{2} + 2\right)^2
            \item \left(\sqrt{2x} + 1\right)^2
            \item \left(3 - \frac{5}{x}\right)^2
            \item \left(2x - \frac{1}{x}\right)^2
            \item (4 - x^2)^2
            \item (x^2 - x)^2
            \item (2x^2 - y)^2
            \item (x^2 + \sqrt{x})^2
            \item (x - 2)^2 (3 - x)^2
            \item \left(\frac{x + 3}{1 - x}\right)^2
            \item (2x + 1)^3
            \item (3 - y)^3
            \item (2 \sqrt[3]{x} - 3)^3        
  \end{mathlist}
  \begin{answers}
      \begin{mathlist}[2]
        \item x^{2} + 4 x + 4
        \item 9 x^{2} + 48 x + 64
        \item x^{4} - 2 \sqrt{5} x^{2} + 5
        \item 4 u^{2} + 28 u v + 49 v^{2}
        \item y^{2} - 8 y + 16
        \item 4 y^{2} - 12 y + 9
        \item x^{2} + 4 x + 4
        \item \frac{x^{2}}{4} + 2 x + 4
        \item 2 \sqrt{2x} + 2 x + 1
        \item 9 - \frac{30}{x} + \frac{25}{x^{2}}
        \item 4 x^{2} - 4 + \frac{1}{x^{2}}
        \item x^{4} - 8 x^{2} + 16
        \item x^{4} - 2 x^{3} + x^{2}
        \item 4 x^{4} - 4 x^{2} y + y^{2}
        \item 2 x^{\frac{5}{2}} + x^{4} + x
        \item x^{4} - 10 x^{3} + 37 x^{2} - 60 x + 36
        \item \frac{x^{2} + 6x + 9}{x^{2} - 2 x + 1}
        \item 8 x^{3} + 12 x^{2} + 6 x + 1
        \item - y^{3} + 9 y^{2} - 27 y + 27
        \item - 36 x^{\frac{2}{3}} + 54 \sqrt[3]{x} + 8 x - 27
    \end{mathlist}
  \end{answers}
\end{exercicio}

    
\begin{exercicio}
  Efetue os produtos abaixo.
  \begin{mathlist}[2]
    \item \left(x + 4\right) \left(x - 4\right)
    \item \left(5 x + 6\right) \left(5 x - 6\right)
    \item \left(2 x + 7 y\right) \left(2 x - 7 y\right)
    \item \left(- x + 2\right) \left(x + 2\right)
    \item \left(\frac{3 x}{2} - \frac{1}{3}\right) \left(\frac{3 x}{2} + \frac{1}{3}\right)
    \item \left(x - \frac{1}{x}\right) \left(x + \frac{1}{x}\right)
    \item \left(y^{2} + 4\right) \left(y^{2} - 4\right)
    \item \left(z - \sqrt{3}\right) \left(z + \sqrt{3}\right)
    \item \left(\sqrt{x} + 5\right) \left(\sqrt{x} - 5\right)
    \item \!\!\left(2 \sqrt{x} \!-\! \sqrt{5}\right)\!\! \left(2 \sqrt{x}\! +\! \sqrt{5}\right)
  \end{mathlist}
  \begin{answers}
      \begin{mathlist}[2]
        \item x^{2} - 16
        \item 25 x^{2} - 36
        \item 4 x^{2} - 49 y^{2}
        \item - x^{2} + 4
        \item \frac{9 x^{2}}{4} - \frac{1}{9}
        \item \frac{x^{4} - 1}{x^{2}}
        \item y^{4} - 16
        \item z^{2} - 3
        \item x - 25
        \item 4 x - 5
    \end{mathlist}
  \end{answers}
\end{exercicio}


               
    \begin{exercicio}
        O número áureo é uma constante real irracional, definida como a raiz positiva da equação quadrática obtida a partir de
        \[ \frac{x + 1}{x} = x \]
        Reescreva a equação acima como uma equação quadrática e determine o número áureo.
    \end{exercicio}

    \begin{exercicio}
        Calcule as expressões abaixo, simplificando o resultado. Sempre que necessário, 
	suponha que os termos no interior das raízes são não negativos e que os denominadores 
	são diferentes de zero.
            \begin{mathlist}
                \item \frac{2 (x - \sqrt{3}) (x + \sqrt{3})}{5 (x^2 - 3)}
                \item \left(\frac{6}{\sqrt{x}} - \frac{4}{\sqrt{y}}\right) \left(\frac{6}{\sqrt{x}} + \frac{4}{\sqrt{y}}\right) \left(\frac{xy}{4}\right)
                \item \frac{(2 - \sqrt{x}) (\sqrt{x} + 2)}{4 - (x - 2)^2}
                \item \frac{(x - \sqrt{2}) (x + \sqrt{2})}{(x - 1)^2 + 2x - 3}
                \item \frac{(\sqrt{x + 5}) (\sqrt{x} - 5)}{\sqrt{x - 25}}
                \item \frac{(\sqrt{2x} + \sqrt{5}) (\sqrt{2x} - \sqrt{5})}{(x - 5)^2 - x^2}
                \item \frac{(x - 3)^2 - (x - \sqrt{3}) (x + \sqrt{3})}{2 - x}
                \item \frac{(x + 2)^2 - (2 \sqrt{x} - 1) (2 \sqrt{x} + 1)}{(5 + x)^2}
            \end{mathlist}     
    \end{exercicio}
                        
    \begin{exercicio}
        Reescreva as expressões abaixo, colocando algum termo em evidência e simplificando 
    	o resultado sempre que possível. Quando necessário, suponha que os denominadores são não nulos.
        \begin{mathlist}[2]
            \item 4 - 2y
            \item 6x - 3
            \item -4x - 10
            \item 35x - 5z + 15y
            \item -10a + 14ab
            \item x^2 - 2x
            \item 8ab - 12b + 4ab^2
            \item 3x^5 - 9x^4 + 18x^7
            \item \frac{3x}{32} - \frac{21}{4}
            \item \frac{5x}{2} - \frac{x^2}{2}
            \item xy + x^2y^2
            \item 4xy + 8yz - 12w^2y
            \item xy^2 + y^5 + 3zy^3
            \item - \frac{5}{12x} + \frac{10}{3x^3}
            \item (4x - 1)^2 + (4x - 1) 3x
            \item (5x + 1) (x - 2) - 4 (x - 2)
            \item \frac{6x^2 - 24x}{3x}
            \item \frac{x (3 - 2x) - 2(3 - 2x)}{x - 2}
        \end{mathlist}
    \end{exercicio}

    
\begin{exercicio}
  Fatore as expressões.
  \begin{mathlist}[2]
    \item x^{2} - 9
    \item 16 x^{2} - 1
    \item 9 - \frac{x^{2}}{4}
    \item x^{2} - 64 y^{2}
    \item 4 y^{2} - 25
    \item 36 x^{2} - 100
    \item 16 - 49 x^{2}
    \item 9 u^{2} - v^{2}
    \item 25 - x^{8}
    \item x^{4} - x^{2}
    \item \frac{9 x^{2}}{4} - \frac{1}{9}
    \item x^{2} - 16
    \item \frac{36}{y^{2}} - \frac{1}{9}
    \item \left(x - 7\right)^{2} - 4
    \item x^{2} + 10 x + 25
    \item 4 x^{2} - 12 x + 9
    \item 3 x^{2} + 12 x + 12
    \item x^{2} - x + \frac{1}{4}
    \item 16 x^{2} + 40 x y + 25 y^{2}
    \item x^{2} y^{2} - 2 x y + 1
    \item x^{2} - 2 \sqrt{3} x + 3
    \item \frac{x^{2}}{4} + \frac{x}{3} + \frac{1}{9}
  \end{mathlist}
  \begin{answers}
      \begin{mathlist}[2]
        \item \left(x - 3\right) \left(x + 3\right)
        \item \left(4 x - 1\right) \left(4 x + 1\right)
        \item - \frac{\left(x - 6\right) \left(x + 6\right)}{4}
        \item \left(x - 8 y\right) \left(x + 8 y\right)
        \item \left(2 y - 5\right) \left(2 y + 5\right)
        \item 4 \left(3 x - 5\right) \left(3 x + 5\right)
        \item - \left(7 x - 4\right) \left(7 x + 4\right)
        \item \left(3 u - v\right) \left(3 u + v\right)
        \item - \left(x^{4} - 5\right) \left(x^{4} + 5\right)
        \item x^{2} \left(x - 1\right) \left(x + 1\right)
        \item \frac{\left(9 x - 2\right) \left(9 x + 2\right)}{36}
        \item \left(x - 4\right) \left(x + 4\right)
        \item - \frac{\left(y - 18\right) \left(y + 18\right)}{9 y^{2}}
        \item \left(x - 9\right) \left(x - 5\right)
        \item \left(x + 5\right)^{2}
        \item \left(2 x - 3\right)^{2}
        \item 3 \left(x + 2\right)^{2}
        \item \frac{\left(2 x - 1\right)^{2}}{4}
        \item \left(4 x + 5 y\right)^{2}
        \item \left(x y - 1\right)^{2}
        \item \left(x-\sqrt{3}\right)^2
	\item \frac{\left(3 x + 2\right)^{2}}{36}
    \end{mathlist}
  \end{answers}
\end{exercicio}

    
\begin{exercicio}
  Resolva as equações abaixo.
  \begin{mathlist}[1]
        \item 4 \left(x - \frac{3}{4}\right) (x - 6) = 0
        \item (x - 9)^2 = 0
        \item (x - 5) (2 - x) = 0
        \item 4x (x + 8) = 0
        \item 8 \left(x + \frac{1}{2}\right) (x - 4) = 0
        \item (5x + 3) (2x + 7) = 0
        \item \left(\frac{x}{4} - \frac{3}{2}\right) \left(3 - \frac{x}{4}\right) = 0
        \item \sqrt{2} \left(x + \sqrt{2}\right) \left(\frac{x}{\sqrt{2}} - 1\right) = 0
  \end{mathlist}
  \begin{answers}
      \begin{mathlist}[2]
        \item x_1 = \sfrac{3}{4} \qquad x_2 = 6
        \item x_1 = 9
        \item x_1 = 2 \qquad x_2 = 5
        \item x_1 = -8 \qquad x_2 = 0
        \item x_1 = - \sfrac{1}{2} \qquad x_2 = 4
        \item x_1 = - \sfrac{7}{2} \qquad x_2 = - \sfrac{3}{5}
        \item x_1 = 6 \qquad x_2 = 12
        \item x_1 = \sqrt{2} \qquad x_2 = - \sqrt{2}
    \end{mathlist}
  \end{answers}
\end{exercicio}

\begin{exercicio}
    Reescreva as equações do exercício anterior na forma $ax^2 + bx + c = 0$.
    \begin{answers}
        \begin{mathlist}[2]
            \item 4 x^{2} - 27 x + 18 = 0 
            \item x^{2} - 18 x + 81 = 0 
            \item - x^{2} + 7 x - 10 = 0 
            \item 4 x^{2} + 32 x = 0 
            \item 8 x^{2} - 28 x - 16 = 0 
            \item 10 x^{2} + 41 x + 21 = 0 
            \item - \frac{x^{2}}{16} + \frac{9 x}{8} - \frac{9}{2} = 0 
            \item x^{2} - 2 = 0 
        \end{mathlist}
    \end{answers}
\end{exercicio}

\begin{exercicio}
    Encontre as soluções reais das equações abaixo, caso existam.
    \begin{mathlist}[2]
        \item x^2 - 10 = 0
        \item 3x^2 - 75 = 0
        \item 4x^2 + 81 = 0
        \item \frac{x^2}{6} - \frac{24}{9} = 0
        \item (x - 2)^2 = 4^2
        \item (2x - 1)^2 - 25 = 0
        \item (x + 3)^2 = \frac{1}{9}
        \item \left(\frac{x}{2} + 1\right)^2 = \frac{9}{4}
    \end{mathlist}
    \begin{answers}
        \begin{mathlist}[2]
            \item x_1 = \sqrt{10} \quad x_2 =\!-\!\sqrt{10}
            \item x_1 = -5 \qquad x_2 = 5 
            \item \) Não há soluções
            \item x_1 = -4 \qquad x_2 = 4 
            \item x_1 = -2 \qquad x_2 = 6 
            \item x_1 = -2 \qquad x_2 = 3 
            \item x_1 = - \sfrac{10}{3} \quad x_2 = -\sfrac{8}{3}
            \item x_1 = -5 \qquad x_2 = 1 
        \end{mathlist}
    \end{answers}
\end{exercicio}

\begin{exercicio}
    Determine as raízes das equações.
    \begin{mathlist}[2]
        \item x^2 - 4x = 0
        \item 5x^2 + x = 0
        \item x^2 = -7x
        \item 2x^2 - 3x = 0
        \item -3x^2 - \frac{x}{2} = 0
        \item \frac{x^2}{3} - \frac{x}{6} = 0
        \item 2x - \sqrt{2}x^2 = 0
        \item \sqrt{3}x - \frac{x^2}{\sqrt{3}} = 0
    \end{mathlist}
    \begin{answers}
        \begin{mathlist}[2]
            \item x_1 = 0 \qquad x_2 = 4 
            \item x_1 = - \sfrac{1}{5} \qquad x_2 = 0 
            \item x_1 = -7 \qquad x_2 = 0 
            \item x_1 = 0 \qquad x_2 = \sfrac{3}{2} 
            \item x_1 = - \sfrac{1}{6} \qquad x_2 = 0 
            \item x_1 = 0 \qquad x_2 = \sfrac{1}{2} 
            \item x_1 = 0 \qquad x_2 = \sqrt{2} 
            \item x_1 = 0 \qquad x_2 = 3 
        \end{mathlist}
    \end{answers}
\end{exercicio}

\begin{exercicio}
    Usando a fórmula de Bhaskara, determine, quando possível, as raízes reais das equações.
    \begin{mathlist}[2]
        \item x^2 - 6x + 8 = 0
        \item x^2 - 2x - 15 = 0
        \item x^2 + 6x + 9 = 0
        \item x^2 + 8x + 12 = 0
        \item 2x^2 + 8x - 10 = 0
        \item x^2 - 6x + 10 = 0
        \item 2x^2 - 7x - 4 = 0
        \item 6x^2 - 5x + 1 = 0 
        \item x^2 - 4x + 13 = 0
        \item 25x^2 - 20x + 4 = 0
        \item x^2 - 2\sqrt{5}x + 5 = 0
        \item 2x^2 - 2\sqrt{2}x - 24 = 0
        \item 3x^2 -0,3x - 0,36 = 0
        \item x^2 - 2,4x + 1,44 = 0
        \item x^2 + 2x + 5 = 0
        \item(x + 8)^2 + 4x = 0
        \item(3 - 4x) (x + 3) = 9
        \item (2 - 3x) (2x - 5) = 4
    \end{mathlist}    
    \begin{answers}
        \begin{mathlist}[2]
            \item x_1 = 2 \qquad x_2 = 4 
            \item x_1 = -3 \qquad x_2 = 5 
            \item x_1 = -3 
            \item x_1 = -6 \qquad x_2 = -2 
            \item x_1 = -5 \qquad x_2 = 1 
            \item \) Não há soluções
            \item x_1 = - \sfrac{1}{2} \qquad x_2 = 4 
            \item x_1 = \sfrac{1}{3} \qquad x_2 = \sfrac{1}{2} 
            \item \) Não há soluções
            \item x_1 = \sfrac{2}{5} 
            \item x_1 = \sqrt{5} 
            \item x_1 = - 2 \sqrt{2} \qquad x_2 = 3 \sqrt{2} 
            \item x_1 = -0,3 \qquad x_2 = 0,4 
            \item \) Não há soluções
            \item \) Não há soluções
            \item x_1 = -16 \qquad x_2 = -4 
            \item x_1 = - \sfrac{9}{4} \qquad x_2 = 0 
            \item x_1 = \sfrac{7}{6} \qquad x_2 = 2 
        \end{mathlist}
    \end{answers}
\end{exercicio}

\begin{exercicio}
    Determine quantas raízes as equações abaixo possuem.
    \begin{mathlist}[2]
        \item 2x^2 + 12x + 18 = 0
        \item x^2 - 3x + 8 = 0
        \item -2x^2 - 5x + 9 = 0
        \item \frac{x^2}{5} - 2x + 20 = 0
        \item -x^2 + 16x - 64 = 0
        \item 3x^2 - 4x + 1 = 0
    \end{mathlist}      
    \begin{answers}
        \begin{mathlist}[2]
            \item n=1 \qquad x_1 = -3 
            \item n=0 
            \item n=2 \quad  x_1 = \frac{-5 \pm\sqrt{97}}{4}
            \item n=0 
            \item n=1 \qquad x_1 = 8 
            \item n=2 \qquad x_1 = \sfrac{1}{3} \quad x_2 = 1 
        \end{mathlist}
    \end{answers}
\end{exercicio}



    \begin{exercicio}
        Determine para que valores de $m$ as equações abaixo possuem ao menos uma raiz.
        \begin{mathlist}[2]
            \item -x^2 - 8x + m = 0
            \item 4x^2 + 12x + m = 0
            \item 5x^2 - 8x + m = 0
            \item mx^2 + 6x - 15 = 0
            \item mx^2 - 5x + 10 = 0
            \item mx^2 - 6x + 9 = 0
        \end{mathlist}
    \end{exercicio}
    
    \begin{exercicio}
        Sabendo que a equação \[4x^2 -(m-3)x + 1 = 0\] possui exatamente uma raiz, $x$, 
        determine os possíveis valores da constante $m$.
    \end{exercicio}

    \begin{exercicio}
        Sabendo que a equação \[mx^2 +(2m+5)x+(m+3) = 0\] não possui raízes reais em $x$, 
        determine os possíveis valores da constante $m$.
    \end{exercicio}
    
    \begin{exercicio}
        A equação $4x^2 -12x + c = 0$ tem duas raízes reais, $x_1$ e $x_2$. 
        Sabendo que $x_2$ é 2 unidades maior que $x_1$, determine essas raízes, bem como o valor de $c$.
    \end{exercicio}
    
    
\begin{exercicio}
  Determine as raízes das equações.
  \begin{mathlist}[2]
    \item 9 x^{4} - 20 x^{2} + 4 = 0
    \item x^{4} + 4 x^{2} - 5 = 0
    \item x^{4} - 8 x^{2} + 16 = 0
    \item x^{4} + 13 x^{2} + 36 = 0
    \item 4 x^{4} - 37 x^{2} + 9 = 0
    \item x^{4} - 24 x^{2} - 25 = 0
  \end{mathlist}
  \begin{answers}
      \begin{mathlist}[1]
        \item x_1 = \sqrt{2} \quad x_2 = - \sqrt{2} \quad x_3 = - \frac{\sqrt{2}}{3} \quad x_4 = \frac{\sqrt{2}}{3} 
        \item x_1 = -1 \quad x_2 = 1 
        \item x_1 = -2 \quad x_2 = 2 
        \item \) Não há soluções
        \item x_1 = -3 \quad x_2 = - \frac{1}{2} \quad x_3 = \frac{1}{2} \quad x_4 = 3 
        \item x_1 = -5 \quad x_2 = 5 
    \end{mathlist}
  \end{answers}
\end{exercicio}


    
    \begin{exercicio}
        Determine o domínio das expressões.
        \begin{mathlist}[2]
            \item \frac{x}{3x - 8}
            \item \frac{y - 12}{16 - y^2}
            \item \frac{\sqrt{3x} + 1}{2x - 15 + x^2}
            \item \frac{2x}{16 + 9x^2}
            \item \sqrt{5x - 4}
            \item \sqrt{35 - 7x}
            \item \sqrt{x^2 - 8}
            \item \sqrt[3]{x - 7}
            \item \frac{\sqrt{2x - 5}}{20 - 8x}
            \item \frac{\sqrt{9 - x^2}}{x - 1}
            \item \frac{x - 6}{\sqrt{x - 5}}
            \item \frac{\sqrt{49 - x^2}}{\sqrt{x}}
            \item \frac{\sqrt{3 - x}}{\sqrt{x - 2}}
        \end{mathlist}
    \end{exercicio}
    
        
\begin{exercicio}
  Simplifique as expressões, fatorando os termos, caso necessário. 
  Suponha sempre que os denominadores são não nulos.
  \begin{mathlist}[2]
      \item \frac{2x - 6}{x - 3}
      \item \frac{2x - 6}{3 - x}
      \item \frac{x^2 - 3x}{4x - 12}
      \item \frac{3y - 12}{6y - 18}
      \item \frac{2x - 4}{3x - 6}
      \item \frac{x^2 - x^3}{x}
      \item \frac{x^2 + x^4}{3x^3}
      \item \frac{x^2y - xy^2}{xy}
      \item \frac{x^2y - xy^2}{x - y}
      \item \frac{x^2 - 9}{x^2 - 3x}
      \item \frac{2x^2 - 50}{x^3 + 5x^2}
      \item \frac{x^2 - 5x + 4}{x - 4}
  \end{mathlist}
  \begin{answers}
      \begin{mathlist}[2]
        \item 2
        \item -2
        \item \frac{x}{4}
        \item \frac{y - 4}{2 \left(y - 3\right)}
        \item \frac{2}{3}
        \item - x \left(x - 1\right)
        \item \frac{x^{2} + 1}{3 x}
        \item x-y
        \item xy
        \item \frac{x + 3}{x}
        \item \frac{2 \left(x - 5\right)}{x^{2}}
        \item x - 1
    \end{mathlist}
  \end{answers}
\end{exercicio}

\begin{exercicio}
  Calcule as expressões abaixo e simplifique o resultado quando possível.
  \begin{mathlist}[2]
      \item \!\!\!\left(\!\frac{x^2 - 4x - 12}{x + 2}\!\right)\!\!\!\left(\!\frac{2x + 1}{x - 6}       \!\right)
      \item \!\!\!\left(\!\frac{2x^2 + 8x}{x - 5}    \!\right)\!\!\!\left(\!\frac{x^2 - 25}{4x^2 + 20x}\!\right)            
      \item \dfrac{\;\dfrac{8}{\;5x\;}\;}{\;\dfrac{4}{\;35x\;}\;}
      \item \dfrac{2 - \dfrac{3}{4}}{\;\dfrac{1}{2x} - \dfrac{1}{3x}\;}
      \item \dfrac{3u^3 v^3}{v^5 u^2} + \dfrac{u^2}{v^2}
      \item \dfrac{2w^3}{y^2} - \dfrac{w^6 y^3}{2w^3 y^5}
      \item \dfrac{2}{x} + \dfrac{4}{5}
      \item \dfrac{2}{5x} - \dfrac{4}{3}
      \item \dfrac{2}{5x - 1} + \dfrac{3}{7}
      \item \dfrac{x + 3}{1 - xd} - 2
      \item \dfrac{5x}{x - 4} + \dfrac{3}{x + 1}
      \item \dfrac{x}{x^2 - 9} - \dfrac{3x - 1}{x - 3}
  \end{mathlist}          
  \begin{answers}
      \begin{mathlist}[2]
        \item 2 x + 1
        \item \frac{x + 4}{2}
        \item 14
        \item \frac{15 x}{2}
        \item \frac{u \left(u + 3\right)}{v^{2}}
        \item \frac{3 w^{3}}{2 y^{2}}
        \item \frac{2 \left(2 x + 5\right)}{5 x}
        \item - \frac{2 \left(10 x - 3\right)}{15 x}
        \item \frac{15 x + 11}{7 \left(5 x - 1\right)}
        \item - \frac{2 d x + x + 1}{d x - 1}
        \item \frac{5 x^{2} + 8 x - 12}{\left(x - 4\right) \left(x + 1\right)}
        \item - \frac{3 x^{2} + 7 x - 3}{\left(x - 3\right) \left(x + 3\right)}
    \end{mathlist}
  \end{answers}
\end{exercicio}

    
    \begin{exercicio}
        Racionalize os denominadores, supondo que $x$ pertença a um domínio adequado.
        \begin{mathlist}[2]
            \item \frac{x}{\sqrt{3x}}
            \item \frac{1}{\sqrt{2 + x}}
            \item \frac{2}{2 - \sqrt{2x}}
            \item \frac{\sqrt{x} + 5}{\sqrt{x} + 4}
            \item \frac{x - 1}{\sqrt{2x - 1} - 1}
            \item \frac{\sqrt{3x}}{\sqrt{x} - \sqrt{3}}
        \end{mathlist}      
    \end{exercicio}
    
    
\begin{exercicio}
  Resolva as equações.
  \begin{mathlist}[2]
      \item \frac{x - 2}{x + 3} = 0
      \item \frac{2x + 5}{x - 1} = 3
      \item \frac{5x - 2}{1 - 3x} = -1
      \item \frac{3 - \sfrac{x}{2}}{3x + 8} = \frac{1}{4}
      \item \frac{3x + 5}{4x - 5} = -3
      \item \frac{4 - \sfrac{x}{2}}{4x + 1} = 0
      \item \frac{x}{x^2 - 3x + 2} = 0
      \item \frac{2x^2}{x + 5} = 5
      \item \frac{x^2}{3x - 2} = 2x - 1
      \item \frac{2 \left(x - \sfrac{5}{6}\right) + 1}{5x - 3} = \frac{2}{3}
      \item \frac{x^2 - 26}{x^2 - 9} = 10
      \item \frac{x^2 + 1}{x^2 - 4} = 6
      \item \frac{x^2 + 3x - 10}{x^2 - 2} = 5
      \item \frac{x^2 - 1}{x^2 - 2x + 1} = 4
      \item \frac{x^2 - 4x + 9}{x^2 - 5x + 6} = 3
      \item \frac{2}{x + 1} - \frac{4}{x - 1} = 0
      \item \frac{4}{x + 1} + \frac{1}{x - 1} = \frac{5}{x^2 - 1}
      \item \frac{3}{x + 1} + \frac{2}{x - 1} = 3
      \item \frac{2}{x - 4} + \frac{5}{x - 2} = 3
      \item \frac{6}{x - 3} + \frac{5}{x - 4} = 2
      \item \frac{x}{x - 2} - \frac{3}{x + 2} = 1
      \item \frac{1}{2x + 1} + \frac{1}{3x - 1} = \frac{2}{5}
      \item \frac{3}{x - 2} - \frac{2}{x + 3} = \frac{1}{x}
      \item \frac{2}{x + 1} - \frac{2}{2x - 3} = \frac{3}{x}
      \item \frac{4}{x - 5} - \frac{2}{x - 2} = \frac{12}{(x - 5) (x - 2)}
      \item \frac{4}{5 - 2x} - \frac{3}{x} = 0
  \end{mathlist}    
  \begin{answers}
      \begin{mathlist}[2]
        \item x_1 = 2 
        \item x_1 = 8 
        \item x_1 = \sfrac{1}{2} 
        \item x_1 = \sfrac{4}{5} 
        \item x_1 = \sfrac{2}{3} 
        \item x_1 = 8 
        \item x_1 = 0 
        \item x_1 = - \sfrac{5}{2} \qquad x_2 = 5 
        \item x_1 = \sfrac{2}{5} \qquad x_2 = 1 
        \item x_1 = 1 
        \item x_1 = - \sfrac{8}{3} \qquad x_2 = \sfrac{8}{3} 
        \item x_1 = \sqrt{5} \qquad x_2 = - \sqrt{5} 
        \item x_1 = 0 \qquad x_2 = \sfrac{3}{4} 
        \item x_1 = \sfrac{5}{3} 
        \item x_1 = 1 \qquad x_2 = \sfrac{9}{2} 
        \item x_1 = -3 
        \item x_1 = \sfrac{8}{5} 
        \item x_1 = - \sfrac{1}{3} \qquad x_2 = 2 
        \item x_1 = 3 \qquad x_2 = \sfrac{16}{3} 
        \item x_1 = \sfrac{7}{2} \qquad x_2 = 9 
        \item x_1 = 10 
        \item x_1 = - \sfrac{1}{12} \qquad x_2 = 2 
        \item x_1 = - \sfrac{1}{2} 
        \item x_1 = - \sfrac{9}{4} \qquad x_2 = 1 
        \item \) Não há soluções
        \item x_1 = \sfrac{3}{2} 
    \end{mathlist}
  \end{answers}
\end{exercicio}

\begin{exercicio}
  Resolva as equações.
  \begin{mathlist}[2]
      \item \frac{3x + 4}{1 - 5x} + \frac{3}{2} = \frac{5}{6}
      \item \sqrt{3x + 4} = 8
      \item \sqrt{x + 1} = 2x - 1
      \item \sqrt{2x + 1} = x - 1
      \item \sqrt{x - 3} + x = 9
      \item \sqrt{4 - x} + 2 = 3x
      \item 4 \sqrt{3x - 1} = \frac{2}{3} - 2x
      \item \sqrt{5 - x^2} = 3 - 2x
      \item \sqrt{8x + 25} - 2 = 3 - 4x
      \item \sqrt{4x + 4} - x + 2 = 0
      \item 2x + \sqrt{10 - 4x} = 1
      \item \sqrt{4x + 5} - x = 2x + 4
      \item \sqrt{2x + 5} - 1 = x
      \item 3 + \sqrt{45 - 6x} = x
      \item \sqrt{2x^2 + 7} = 2x - 1
      \item \sqrt{25 - 3x^2} = -x
      \item 2 \sqrt{9x^2 - 7} = 6
      \item \sqrt{x^2 + 3} + x = 5
      \item \sqrt{4x^2 - 1} + 1 = 2x
      \item \sqrt{9x^2 - 2} - 3x = 4
      \item (x + 2)^{2/3} = 9
      \item (5x - 6)^{3/2} = 8
      \item (4x^2 - 13)^{3/5} = 27
      \item x^{1/2} - 13x^{1/4} + 12 = 0
      \item 3 \sqrt{x} - \sqrt[4]{x} = 2
      \item x^{1/3} - x^{1/6} = 2
  \end{mathlist}
  \begin{answers}
      \begin{mathlist}[2]
        \item x_1 = 14 
        \item x_1 = 20 
        \item x_1 = \sfrac{5}{4} 
        \item x_1 = 4 
        \item x_1 = 7 
        \item x_1 = \sfrac{11}{9} 
        \item x_1 = \sfrac{1}{3} 
        \item x_1 = \sfrac{2}{5} 
        \item x_1 = 0 
        \item x_1 = 8 
        \item x_1 = - \sfrac{3}{2} 
        \item x_1 = - \sfrac{11}{9} \qquad x_2 = -1 
        \item x_1 = 2 
        \item x_1 = 6 
        \item x_1 = 3 
        \item x_1 = - \sfrac{5}{2} 
        \item x_1 = - \sfrac{4}{3} \qquad x_2 = \sfrac{4}{3} 
        \item x_1 = \sfrac{11}{5} 
        \item x_1 = \sfrac{1}{2} 
        \item x_1 = - \sfrac{3}{4} 
        \item x_1 = 25 
        \item \) Não há soluções
        \item \) Não há soluções
        \item x_1 = 1 \qquad x_2 = 20736 
        \item x_1 = 1 
        \item x_1 = 64 
    \end{mathlist}
  \end{answers}
\end{exercicio}


    

\begin{exercicio}
  Efetue as operações indicadas e simplifique cada expressão.
  \begin{mathlist}
      \item (7x^2 - 2x + 5) + (2x^2 + 5x - 4)
      \item (3x^2 + 5xy + 2y) + (4 - 3xy - 2x^2)
      \item x - \{2x - [-x - (1 - x)]\}
      \item \left(\frac{1}{3} - 1 + e\right) - \left(- \frac{1}{3} - 1 + e^{-1}\right)
      \item - \frac{3}{4} y - \frac{1}{4} x + 100 + \frac{1}{2} x + \frac{1}{4} y - 120
      \item 3 \sqrt{8} + 8 - 2 \sqrt{y} + \frac{1}{2} \sqrt{x} - \frac{3}{4} \sqrt{y}    
      \item \frac{8}{9} x^2 + \frac{2}{3} x + \frac{16}{3} x^2 - \frac{16}{3} x - 2x + 2
      \item (x + 8) (x - 2)
      \item (3a - 4b)^2
      \item (x + 2y)^2
      \item (2x + y) (2x - y)
      \item (x^{1/2} + 1) \left(\frac{1}{2} x^{-1/2}\right) - (x^{1/2} - 1) \left(\frac{1}{2} x^{-1/2}\right)
      \item 2 (t + \sqrt{t})^2 - 2t^2
  \end{mathlist}
  \begin{answers}
      \begin{mathlist}[2]
        \item 9 x^{2} + 3 x + 1
        \item x^{2} + 2 x y + 2 y + 4
        \item - x - 1
        \item \frac{2}{3} + e - \frac{1}{e}
        \item \frac{x}{4} - \frac{y}{2} - 20
        \item \frac{\sqrt{x}}{2} - \frac{11 \sqrt{y}}{4} + 8 + 6 \sqrt{2}
        \item \frac{56 x^{2}}{9} - \frac{20 x}{3} + 2
        \item x^{2} + 6 x - 16
        \item 9 a^{2} - 24 a b + 16 b^{2}
        \item x^{2} + 4 x y + 4 y^{2}
        \item 4 x^{2} - y^{2}
        \item \frac{1}{\sqrt{x}}
        \item 4 t^{\frac{3}{2}} + 2 t
    \end{mathlist}
  \end{answers}
\end{exercicio}


\begin{exercicio}
  Coloque em evidência o maior fator comum de cada expressão.
  \begin{mathlist}[2]
      \item 4x^5 - 12x^4 - 6x^3
      \item 4x^2 y^2 z - 2x^5 y^2 + 6x^3 y^2 z^2
      \item 3x^{2/3} - 2x^{1/3}
      \item e^{-x} - xe^{-x}
      \item 2ye^{xy^2} + 2xy^3 e^{xy^2}
      \item 2x^{-5/2} - \frac{3}{2} x^{-3/2}
      \item \frac{1}{2} \left(\frac{2}{3} u^{3/2} - 2u^{1/2}\right)
  \end{mathlist}
  \begin{answers}
      \begin{mathlist}[2]
        \item 2 x^{3} \left(2 x^{2} - 6 x - 3\right)
        \item - 2 x^{2} y^{2} \left(x^{3} - 3 x z^{2} - 2 z\right)
        \item \sqrt[3]{x} \left(3 \sqrt[3]{x} - 2\right)
        \item -e^{-x} \left(x - 1\right) 
        \item 2 ye^{x y^{2}} \left(x y^{2} + 1\right) 
        \item \frac{-3 + \frac{4}{x}}{2 x^{\frac{3}{2}}}
        \item \frac{\sqrt{u} \left(u - 3\right)}{3}
    \end{mathlist}
  \end{answers}
\end{exercicio}


\begin{exercicio}
  Fatore cada expressão.
  \begin{mathlist}[2]
      \item 6ac + 3bc - 4ad - 2bd
      \item 4a^2 - b^2
      \item 12x^2 - 3y^2
      \item 10 - 14x - 12x^2
      \item 3x^2 - 6x - 24
      \item 12x^2 - 2x - 30
      \item (x + y)^2 - 1
      \item 8a^2 - 2ab - 6b^2
      \item x^6 + 125
  \end{mathlist}
  \begin{answers}
      \begin{mathlist}[2]
        \item \left(2 a + b\right) \left(3 c - 2 d\right)
        \item \left(2 a - b\right) \left(2 a + b\right)
        \item 3 \left(2 x - y\right) \left(2 x + y\right)
        \item - 2 \left(2 x - 1\right) \left(3 x + 5\right)
        \item 3 \left(x - 4\right) \left(x + 2\right)
        \item 2 \left(2 x + 3\right) \left(3 x - 5\right)
        \item \left(x + y - 1\right) \left(x + y + 1\right)
        \item 2 \left(a - b\right) \left(4 a + 3 b\right)
        \item \left(x^{2} + 5\right) \left(x^{4} - 5 x^{2} + 25\right)
    \end{mathlist}
  \end{answers}
\end{exercicio}


\begin{exercicio}
  Efetue as operações indicadas e simplifique cada expressão.
  \begin{mathlist}
      \item (x^2 + y^2) x - xy (2y)
      \item 2kr (R - r) - kr^2
      \item 2 (x - 1) (2x + 2)^3 [4(x - 1) + (2x + 2)]
      \item 5x^2 (3x^2 + 1)^4 (6x) + (3x^2 + 1)^5 (2x)
      \item 4 (x - 1)^2 (2x + 2)^3 (2) + (2x + 2)^4 (2) (x - 1)
      \item (x^2 + 2)^2 [5 (x^2 + 2)^2 - 3] (2x)
      \item (x^2 - 4) (x^2 + 4) (2x + 8) - (x^2 + 8x - 4) (4x^3)
  \end{mathlist}
  \begin{answers}
      \begin{mathlist}[1]
        \item x \left(x - y\right) \left(x + y\right)
        \item - k r \left(- 2 R + 3 r\right)
        \item 32 \left(x - 1\right) \left(x + 1\right)^{3} \left(3 x - 1\right)
        \item 2 x \left(3 x^{2} + 1\right)^{4} \left(18 x^{2} + 1\right)
        \item 32 \left(x - 1\right) \left(x + 1\right)^{3} \left(3 x - 1\right)
        \item 2 x \left(x^{2} + 2\right)^{2} \left(5 x^{4} + 20 x^{2} + 17\right)
        \item - 2 \left(x^{5} + 12 x^{4} - 8 x^{3} + 16 x + 64\right)
    \end{mathlist}
  \end{answers}
\end{exercicio}

    
\end{exercicios}

%------------------------------------------------------------------------------%

%------------------------------------------------------------------------------%
% \chapter{Revisão}
%------------------------------------------------------------------------------%

%------------------------------------------------------------------------------%
\begin{exercicios}{Polinômios}
    
    \begin{exercicio}
        Defina uma função $f(x)$ que forneça a área da região destacada na figura, lembrando que a área de um retângulo de lados $b$ e $h$ é $bh$.
        \addgraphic{.8}{exercicio-polinomio-area-retangulo}
    \end{exercicio}
    
    \begin{exercicio}
        Encontre as raízes de cada equação por fatoração.
        \begin{mathlist}[2]
            \item x^2 + x - 12 = 0
            \item 3x^2 - x - 4 = 0
            \item \frac{1}{4} x^2 - x + 1 = 0
            \item 8x^2 + 2x - 3 = 0
            \item -6x^2 - 10x + 4 = 0
            \item 2x^4 + x^2 = 1
        \end{mathlist}
        \begin{answers}
            \begin{mathlist}[2]
                \item x_1 = -4 \qquad x_2 = 3
                \item x_1 = -1 \qquad x_2 = \sfrac{4}{3}
                \item x_1 = 2
            \end{mathlist}
        \end{answers}
    \end{exercicio}
    
    \begin{exercicio}
        Resolva as equações utilizando a fórmula quadrática.
        \begin{mathlist}[2]
            \item 4x^2 + 5x - 6 = 0
            \item 8x^2 - 8x - 3 = 0
            \item x^2 - 6x + 6 = 0
            \item 2x^2 + 8x + 7 = 0
        \end{mathlist}
        \begin{answers}
            \begin{mathlist}[2]
                \item x_1 = -2 \qquad x_2 = \sfrac{4}{3}
                \item x_{1,2} = \frac{2\pm\sqrt{10}}{4}
                \item x_{1,2} = 3\pm\sqrt{3}
                \item x_{1,2} = -2\pm\frac{\sqrt{2}}{2}
            \end{mathlist}
        \end{answers}    
    \end{exercicio}
    
    \begin{exercicio}
        Encontre as raízes da equação $2x^4 + x^2 = 1$ substituindo $x^2$ por $y$.
        \begin{answers}
            \(\qquad x_{1,2} = \pm\frac{1}{\sqrt{2}} \)
        \end{answers}
    \end{exercicio}    
    
    \begin{exercicio}
        Resolva as inequações quadráticas.
        \begin{mathlist}[2]
            \item   x^2 + 3x \geq 10
            \item -3x^2 - 11x + 4 > 0
            \item -4x^2 + 4x - 1 < 0
            \item   x^2 + x + 2 \leq 0
        \end{mathlist}
        \begin{answers}
            \begin{mathlist}[2]
                \item x \leq -5 \text{ ou } x \geq 2
                \item -4 < x < 1/3
                \item x \neq 1/2
                \item \) Não há solução.
            \end{mathlist}
        \end{answers}
    \end{exercicio}
    
    \begin{exercicio}
        Para cada expressão na forma $p(x)/d(x)$ abaixo, calcule o quociente $q(x)$ e o resto $r(x)$.
        \begin{mathlist}
            \item (2x^3 - 3x^2 + 6)/(x^2 - 2)
            \item (6x^2 - 4x - 3)/(3x - 5)
            \item (x^4 + 2x - 12)/(x + 2)
            \item (4x^3 + 2x^2 + 11x)/(2x^2 + 3)
            \item (6x^4 + 5x^3 - 2x)/(3x - 2)
            \item (4x^3 + 6x - 10)/(2x - 4)
            \item (x^2 - 5x + 8)/(x - 3)
            \item (3x + 7)/(x + 4)
            \item (x^4 - 2)/(x - 1)
            \item (24x^3 - 4x - 1)/(2x - 1)
            \item (8x^3 - 12x^2 - 2x)/(4x - 8)
            \item (x^3 - 3x^2 + 4x - 5)/(x - 4)
            \item (2x^4 - 4x^3 + x - 17)/(x^2 - 4)
            \item (x^4 - 6x^3 + 3x^2 - 2x + 3)/(x^2 - 2x - 3)
            \item (x^4 - 5x^2 + 4)/(x^2 - 1)
            \item (3x^5 - 2x^3 - 11x)/(x^3 - 3x)
            \item (6x^2 + 7x + 9)/(2x^2 - 5x + 1)
        \end{mathlist}
        \begin{answers}
            \begin{alphalist}
                \item $q(x) = 2x - 3                $, $r(x) = 4x$
                \item $q(x) = 2x + 2                $, $r(x) = 7$
                \item $q(x) = x^3 - 2x^2 + 4x - 6   $, $r(x) = 0$
                \item $q(x) = 2x + 1                $, $r(x) = 5x - 3$
                \item $q(x) = 2x^3 + 3x^2 + 2x + 2/3$, $r(x) = 4/3$
                \item $q(x) = 2x^2 + 4x + 11        $, $r(x) = 34$
                \item $q(x) = x - 2                 $, $r(x) = 2$
                \item $q(x) = 3                     $, $r(x) = -5$
                \item $q(x) = x^3 + x^2 + x + 1     $, $r(x) = -1$
                \item $q(x) = 12x^2 + 6x + 1        $, $r(x) = 0$
                \item $q(x) = 2x^2 + x + 3/2        $, $r(x) = 12$
                \item $q(x) = x^2 + x + 8           $, $r(x) = 27$
                \item $q(x) = 2x^2 - 4x + 8         $, $r(x) = 15 - 15x$
                \item $q(x) = x^2 - 4x - 2          $, $r(x) = -18x - 3$
                \item $q(x) = x^2 - 4               $, $r(x) = 0$
                \item $q(x) = 3x^2 + 7              $, $r(x) = 10x$
                \item $q(x) = 3                     $, $r(x) = 22x + 6$
            \end{alphalist}
        \end{answers}
    \end{exercicio}
    
    \begin{exercicio}
        Para os problemas do exercício anterior, expresse $p(x)$ na forma $d(x)q(x) + r(x)$.
        \begin{answers}
            \begin{mathlist}
                \item p(x) = (x^2 - 2)(2x - 3) + 4x
                \item p(x) = (3x - 5)(2x + 2) + 7
                \item p(x) = (x + 2)(x^3 - 2x^2 + 4x - 6)
                \item p(x) = (2x^2 + 3)(2x + 1) + 5x - 3
                \item p(x) = (3x-2)(2x^3 +3x^2 +2x+2/3)+4/3
                \item p(x) = (2x - 4)(2x^2 + 4x + 11) + 34
                \item p(x) = (x - 3)(x - 2) + 2
                \item p(x) = (x + 4)3 - 5
                \item p(x) = (x - 1)(x^3 + x^2 + x + 1) - 1
                \item p(x) = (2x - 1)(12x^2 + 6x + 1)
                \item p(x) = (4x - 8)(2x^2 + x + 3/2 ) + 12
                \item p(x) = (x - 4)(x^2 + x + 8) + 27
                \item p(x) = (x^2 - 4)(2x^2 - 4x + 8) - 15x + 15
                \item p(x) = (x^2 -2x-3)(x^2 -4x-2)-18x-3
                \item p(x) = (x^2 - 1)(x^2 - 4)
                \item p(x) = (x^3 - 3x)(3x^2 + 7) + 10x
                \item p(x) = (2x^2 - 5x + 1)3 + 22x + 6
            \end{mathlist}
        \end{answers}
    \end{exercicio}
    
    \begin{exercicio}
        Para cada expressão na forma $p(x)/d(x)$ abaixo, calcule o quociente $q(x)$ e o resto $r(x)$ usando o algoritmo de Ruffini.
        \begin{mathlist}
            \item (x^4 + 2x - 12)/(x + 2)
            \item (3x^2 + 2x - 5)/(x - 2)
            \item (4x^4 + 6x^3 - 8x^2 + 22x - 24)/(x + 3)
            \item (-2x^3 + 3x^2 + 12x + 25)/(x - 4)
            \item (x^5 - 9x^3 + 2x)/(x - 3)
            \item (-6x^3 + 4x^2 - x + 2)/(x - 1/3)
            \item (2x^3 - 9x^2 + 6x + 5)/(x - 3/2)
            \item (x^2 - 5x - 6)/(x + 1)
            \item (-4x^2 + 11x + 26)/(x - 4)
            \item (6x^2 - 7x - 9)/(x + 1/2)
            \item (x^3 - 9x^2 )/(x - 3)
            \item (5x^4 - 1)/(x - 2)
            \item (8x^4 + 6x^3 + 3x^2 + 1)/(x - 1/2)
            \item (x^4 - 20x^2 - 50)/(x - 5)
            \item x^4 /(x - 3)
            \item (2x^5 - 4x^4 + 9x^3 - 5x^2 + x - 3)/(x - 1)
        \end{mathlist}
        \begin{answers}
            \begin{alphalist}
                \item $q(x) = x^3 - 2x^2 + 4x - 6$,        $r(x) = 0$
                \item $q(x) = 3x + 8$,                     $r(x) = 11$
                \item $q(x) = 4x^3 - 6x^2 + 10x - 8$,      $r(x) = 0$
                \item $q(x) = -2x^2 - 5x - 8$,             $r(x) = -7$
                \item $q(x) = x^4 + 3x^3 + 2$,             $r(x) = 6$
                \item $q(x) = -6x^2 + 2x - 1/3$,           $r(x) = 17/9$
                \item $q(x) = 2x^2 - 6x - 3$,              $r(x) = 1/2$
                \item $q(x) = x - 6$,                      $r(x) = 0$
                \item $q(x) = -4x - 5$,                    $r(x) = 6$
                \item $q(x) = 6x - 10$,                    $r(x) = -4$
                \item $q(x) = x^2 - 6x - 18$,              $r(x) = -54$
                \item $q(x) = 5x^3 + 10x^2 + 20x + 40$,    $r(x) = 79$
                \item $q(x) = 8x^3 + 10x^2 + 8x + 4$,      $r(x) = 3$
                \item $q(x) = x^3 + 5x^2 + 5x + 25$,       $r(x) = 75$
                \item $q(x) = x^3 + 3x^2 + 9x + 27$,       $r(x) = 81$
                \item $q(x) = 2x^4 -2x^3 +7x^2 +2x+3$,     $r(x) = 0$
            \end{alphalist}
        \end{answers}
    \end{exercicio}
    
    \begin{exercicio}
        Para os problemas do exercício anterior, expresse $p(x)/d(x)$ na forma $q(x) + r(x)/d(x)$.
        \begin{answers}
            \begin{mathlist}
                \item p(x)/d(x) = x^3 - 2x^2 + 4x - 6
                \item p(x)/d(x) = 3x + 8 + 11/(x - 2)
                \item p(x)/d(x) = 4x^3 - 6x^2 + 10x - 8
                \item p(x)/d(x) = -2x^2 - 5x - 8 - 7/(x - 4)
                \item p(x)/d(x) = x^4 + 3x^3 + 2 + 6/(x - 3)
                \item p(x)/d(x) = -6x^2 + 2x -1/3+ 17/(9x - 3)
                \item p(x)/d(x) = 2x^2 - 6x - 3 + 1/(2x - 3)
                \item p(x)/d(x) = x - 6
                \item p(x)/d(x) = -4x - 5 + 6/(x - 4)
                \item p(x)/d(x) = 6x - 10 - 4/(x + 12 )
                \item p(x)/d(x) = x^2 - 6x - 18 - 54/(x - 3)
                \item p(x)/d(x) = 5x^3 + 10x^2 + 20x + 40 + 79/(x - 2)
                \item p(x)/d(x) = 8x^3 +10x^2 +8x+4+3/(x- 1/2 )
                \item p(x)/d(x) = x^3 +5x^2 +5x+25+75/(x-5)
                \item p(x)/d(x) = x^3 +3x^2 +9x+27+81/(x-3)
                \item p(x)/d(x) = 2x^4 - 2x^3 + 7x 2 + 2x + 3
            \end{mathlist}
        \end{answers}
    \end{exercicio}
    
    \begin{exercicio}
        Usando o algoritmo de Ruffini, verifique quais valores abaixo correspondem a zeros das funções associadas.
        \begin{alphalist}
            \item $f(x) = x^2 - 3x + 4$,                         $x_1 = 2$,    $x_2 = -2$
            \item $f(x) = -2x^2 + 3x + 2$,                       $x_1 = -1/2$, $x_2 = -2$
            \item $f(x) = 4x + x^2 $,                            $x_1 = -4$,   $x_2 = 0$
            \item $f(x) = -x^2 - 4$,                             $x_1 = 2$,    $x_2 = -2$
            \item $f(x) = x^3 - 4x^2 + 6x - 9$,                  $x_1 = 3$,    $x_2 = 1$
            \item $f(x) = x^4 - 3x^2 + 2$,                       $x_1 = 5$,    $x_2 = -1$
            \item $f(x) = 2x^4 + x^3 - 25x^2 + 12x$,             $x_1 = -4$,   $x_2 = 3$
            \item $f(x) = x^5 + 2x^4 - 3x^3 + 12x^2 - 28x + 16$, $x_1 = 6$,    $x_2 = 2$
            \item $f(x) = 9x^3 - 15x^2 - 26x + 40$,              $x_1 = 1/2$,  $x_2 = 4/3$
            \item $f(x) = x^3 - 21x - 20$,                       $x_1 = 5$,    $x_2 = 1$
        \end{alphalist}
        \begin{answers}
            \begin{alphalist}
                \item Nenhum valor é um zero da função.
                \item Só -1/2 é um zero de $f$
                \item -4 e 0 são zeros de $f$
                \item Nenhum valor é um zero da função.
                \item Só 3 é um zero de $f$
                \item Só -1 é um zero de $f$
                \item -4 e 3 são zeros de $f$
                \item Nenhum valor é um zero da função.
                \item Só 4/3 é um zero de $f$
                \item Só 5 é um zero de $f$
            \end{alphalist}
        \end{answers}
    \end{exercicio}
       
    \begin{exercicio}
        Para cada função polinomial abaixo, determine o valor da constante $c$ de modo que o termo fornecido seja um fator de $p$.
        \begin{alphalist}
            \item $p(x) = x^2 - 9x + c$,                   \tabto{55mm} fator: $x-8$
            \item $p(x) = 5x^2 + cx + 9$,                  \tabto{55mm} fator: $x+3$
            \item $p(x) = x^3 - 6x^2 + 3x + c$,            \tabto{55mm} fator: $x-5$
            \item $p(x) = 3x^3 + cx^2 - 13x + 3$,          \tabto{55mm} fator: $x-1$
            \item $p(x) = x^4 - 2x^3 + 8x^2 + cx - 2$,     \tabto{55mm} fator: $x-2$
            \item $p(x) = 2x^4 - 10x^3 + cx^2 + 6x + 40$,  \tabto{55mm} fator: $x-4$
        \end{alphalist}
        \begin{answers}
            \begin{mathlist}[3]
                \item c = 8
                \item c = 18
                \item c = 10
                \item c = 7
                \item c = -15
                \item c = 4
            \end{mathlist}
        \end{answers}
    \end{exercicio}
    
    \begin{exercicio}
        Determine as raízes das equações abaixo. Escreva na forma fatorada os polinômios que aparecem no lado esquerdo das equações.
        \begin{mathlist}[2]
            \item   x^3 -  4x           = 0
            \item   x^3 -  4x^2 - 21x   = 0
            \item  2x^3 + 11x^2 -  6x   = 0
            \item -3x^3 +  6x^2 +  9x   = 0
            \item   x^4 -   x^3 - 20x^2 = 0
            \item   x^4 -  8x^3 + 16x^2 = 0
            \item  5x^4 -  8x^3 +  3x^2 = 0
            \item  8x^4 -  6x^3 -  2x^2 = 0
        \end{mathlist}
        \begin{answers}
            \begin{alphalist}
                \item Raízes: 0, 2 e -2,                 \\[1mm] $p(x) = x(x - 2)(x + 2)$
                \item Raízes: 0, -3 e 7,                 \\[1mm] $p(x) = x(x + 3)(x - 7)$
                \item Raízes: 0, 1/2 e -6,               \\[1mm] $p(x) = 2x (x - 1/2 ) (x + 6)$
                \item Raízes: 0, 3 e -1,                 \\[1mm] $p(x) = -3x(x - 3)(x + 1)$
                \item Raízes: 5, -4 e 0 (mult. 2),       \\[1mm] $p(x) =  x^2 (x - 5)(x + 4)$
                \item Raízes: 4 (mult. 2) e 0 (mult. 2), \\[1mm] $p(x) =  x^2 (x - 4)^2$
                \item Raízes: 3/5, 1 e 0 (mult. 2),      \\[1mm] $p(x) = 5x^2 (x - 3/5 ) (x - 1)$
                \item Raízes: 3/2, -1/4 e 0 (mult. 2),   \\[1mm] $p(x) = 8x^2 (x - 3/2 ) (x + 1/4)$
            \end{alphalist}
        \end{answers}
    \end{exercicio}
    
    \begin{exercicio}
        Determine as raízes das equações abaixo. Escreva na forma fatorada os polinômios que aparecem no lado esquerdo das equações.
        \begin{alphalist}
            \item $ x^3 +   x^2 -  2x   - 2         = 0$, \\[1mm] sabendo que $x =  -1$              é uma raiz.
            \item $ x^3 -  5x^2 -  4x   + 20        = 0$, \\[1mm] sabendo que $x =   2$              é uma raiz.
            \item $ x^4 -  9x^3 -   x^2 + 81x - 72  = 0$, \\[1mm] sabendo que $x =   8$ e $x = 3$    são raízes.
            \item $ x^3 -  3x^2 - 10x   + 24        = 0$, \\[1mm] sabendo que $x =   4$              é uma raiz.
            \item $ x^3 -  4x^2 - 17x   + 60        = 0$, \\[1mm] sabendo que $x =   3$              é uma raiz.
            \item $4x^4 - 21x^3 - 19x^2 + 6x        = 0$, \\[1mm] sabendo que $x = 1/4$              é uma raiz.
            \item $4x^3 - 16x^2 + 21x   - 9         = 0$, \\[1mm] sabendo que $x =   1$              é uma raiz.
            \item $3x^3 - 26x^2 + 33x   + 14        = 0$, \\[1mm] sabendo que $x =   7$              é uma raiz.
            \item $ x^4 -  9x^3 + 17x^2 + 33x - 90  = 0$, \\[1mm] sabendo que $x =  -2$ e $x = 5$    são raízes.
            \item $ x^4 -  6x^3 -  5x^2 + 30x       = 0$, \\[1mm] sabendo que $x =   6$              é uma raiz.
            \item $2x^4 +  9x^3 - 80x^2 + 21x + 108 = 0$, \\[1mm] sabendo que $x =   4$ e $x = 3/2$  são raízes.
            \item $ x^3 +  7x^2 + 13x   + 15        = 0$, \\[1mm] sabendo que $x =  -5$              é uma raiz.
            \item $3x^3 +  2x^2 + 17x   - 6         = 0$, \\[1mm] sabendo que $x = 1/3$              é uma raiz.
            \item $ x^3 +  7x^2 + 20x   + 32        = 0$, \\[1mm] sabendo que $x =  -4$              é uma raiz.
            \item $ x^3 -  3x^2 +  9x   - 27        = 0$, \\[1mm] sabendo que $x =   3$              é uma raiz.
            \item $2x^3 - 10x^2 - 13x   - 105       = 0$, \\[1mm] sabendo que $x =   7$              é uma raiz.
            \item $ x^4 +  2x^3 -  5x^2 - 36x + 60  = 0$, \\[1mm] sabendo que $x =   2$              é uma raiz de multiplicidade 2.
            \item $ x^4 -  6x^3 + 25x^2 - 150x      = 0$, \\[1mm] sabendo que $x =   6$              é uma raiz.
            \item $6x^4 +  7x^3 +  6x^2 - 1         = 0$, \\[1mm] sabendo que $x = 1/3$ e $x =-1/2$  são raízes.
        \end{alphalist}
        \begin{answers}
            \begin{alphalist}
                \item $p(x) = (x + 1)(x - \sqrt{2})(x + \sqrt{2})$,   \\[1mm] Raízes: -1, $\sqrt{2}$ e $-\sqrt{2}$        
                \item $p(x) = (x - 1)(x - 2)(x - 8)$,                 \\[1mm] Raízes: 1, 2 e 8
                \item $p(x) = (x + 3)(x - 1)(x - 3)(x - 8)$,          \\[1mm] Raízes: -3, 1, 3 e 8             
                \item $p(x) = (x + 3)(x - 2)(x - 4)$,                 \\[1mm] Raízes: -3, 2 e 4
                \item $p(x) = (x + 4)(x - 3)(x - 5)$,                 \\[1mm] Raízes: -4, 3 e 5
                \item $p(x) = 4(x + 1/4 )x(x + 1)(x - 6)$,            \\[1mm] Raízes: -1, 0, 1/4 e 6
                \item $p(x) = 4 (x - 3/2 ) (x - 1)$,                  \\[1mm] Raízes: 1 e 3/2
                \item $p(x) = (x + 5) (x 2 + 2x + 3)$,                \\[1mm] Raiz: -5
                \item $p(x) = 3 (x - 1/3)(x^2+x+6)$,                  \\[1mm] Raiz: 1/3
                \item $p(x) = (x + 4) (x^2 + 3x + 8)$,                \\[1mm] Raiz: -4
                \item $p(x) = (x - 3) (x^2 + 9)$,                     \\[1mm] Raiz: 3
                \item $p(x) = (x - 7) (2x^2 + 4x + 15)$,              \\[1mm] Raiz: 7
                \item $p(x) = (x - 2)^2 (x^2 + 6x + 15)$,             \\[1mm] Raiz: 2 (multiplicidade 2)
                \item $p(x) = x(x - 6) (x^2 + 25)$,                   \\[1mm] Raizes: 0 e 6
                \item $p(x) = 6 (x + 1/2 ) (x - 1/3 ) (x^2 + x + 1)$, \\[1mm] Raizes: -1/2 e 1/3                            
            \end{alphalist}
        \end{answers}
    \end{exercicio}
    
\end{exercicios}

%------------------------------------------------------------------------------%

%------------------------------------------------------------------------------%
% \chapter{Revisão}
%------------------------------------------------------------------------------%

%------------------------------------------------------------------------------%
\begin{exercicios}{Expressões Racionais}
    
    
\begin{exercicio}
  Simplifique as expressões.
  \begin{mathlist}[2]
      \item \frac{x^2 + x - 2}{x^2 - 4}
      \item \frac{2a^2 - 3ab - 9b^2}{2ab^2 + 3b^3}
      \item \frac{12t^2 + 12t + 3}{4t^2 - 1}
      \item \frac{3(4x - 1) - 4(3x + 1)}{(4x - 1)^2}
  \end{mathlist}    
  \begin{answers}
      \begin{mathlist}[2]
        \item \frac{x - 1}{x - 2}
        \item \frac{a - 3 b}{b^{2}}
        \item \frac{3 \left(2 t + 1\right)}{2 t - 1}
        \item - \frac{7}{\left(4 x - 1\right)^{2}}
    \end{mathlist}
  \end{answers}
\end{exercicio}

\begin{exercicio}
  Efetue as operações indicadas e simplifique cada expressão.
  \begin{mathlist}
      \item \frac{2a^2 - 2b^2}{b - a} \, \frac{4a + 4b}{a^2 + 2ab + b^2}
      \item \frac{x^2 - 6x + 9}{x^2 - x - 6} \, \frac{3x + 6}{2x^2 - 7x + 3}
      \item \frac{3x^2 + 2x - 1}{2x + 6} \div \frac{x^2 - 1}{x^2 + 2x - 3}
      \item \frac{58}{3 (3t + 2)} + \frac{1}{3}
      \item \frac{2x}{2x - 1} - \frac{3x}{2x + 5}
      \item \frac{-xe^x}{x + 1} + e^x
      \item \frac{4}{x^2 - 9} - \frac{5}{x^2 - 6x + 9}
      \item \frac{x}{1 - x} + \frac{2x + 3}{x^2 - 1}
      \item \frac{1 + \frac{1}{x}}{1 - \frac{1}{x}}
      \item \frac{4x^2}{\sqrt{2x^2 + 7}} + \sqrt{2x^2 + 7}
      \item 6 (2x + 1)^2 \sqrt{x^2 + x} + \frac{(2x +1)^4}{2 \sqrt{x^2 + x}}
      \item \frac{2x (x + 1)^{-1/2} - (x + 1)^{1/2}}{x^2}
      \item \frac{(x^2 + 1)^{1/2} - 2x^2 (x^2 + 1)^{-1/2}}{1 - x^2}
      \item \frac{(2x + 1)^{1/2} - (x + 2) (2x + 1)^{-1/2}}{2x + 1}
  \end{mathlist}        
  \begin{answers}
      \begin{mathlist}[2]
        \item -8
        \item \frac{3}{2 x - 1}
        \item \frac{3 x}{2} - \frac{1}{2}
        \item \frac{t + 20}{3 t + 2}
        \item \frac{x \left(13 - 2x\right)}{\left(2 x - 1\right) \left(2 x + 5\right)}
        \item \frac{e^{x}}{x + 1}
        \item - \frac{x + 27}{x^{3} - 3 x^{2} - 9 x + 27}
        \item \frac{- x^{2} + x + 3}{x^{2} - 1}
        \item \frac{x + 1}{x - 1}
        \item \frac{7+6\,x^2}{\sqrt{2\,x^2+7}}
        \item \frac{\left(1+2x\right)^2\left(1+16x+16x^2\right)}{2\sqrt{x\left(x+1\right)}}
        \item \frac{x - 1}{x^{2} \sqrt{x + 1}}
        \item \frac{1}{\sqrt{x^{2} + 1}}
        \item \frac{x - 1}{\left(2 x + 1\right)^{\frac{3}{2}}}
    \end{mathlist}
  \end{answers}
\end{exercicio}


    
%    \begin{exercicio}
%        Simplifique as expressões.
%        \begin{mathlist}
%            \item \frac{x^2 + x - 2}{x^2 - 4}
%            \item \frac{2a^2 - 3ab - 9b^2}{2ab^2 + 3b^3}
%            \item \frac{12t^2 + 12t + 3}{4t^2 - 1}
%            \item \frac{(4x - 1) (3) - (3x + 1) (4)}{(4x - 1)^2}
%        \end{mathlist}    
%    \end{exercicio}
%    
%    \begin{exercicio}
%        Efetue as operações indicadas e simplifique cada expressão.
%        \begin{mathlist}
%            \item \frac{2a^2 - 2b^2}{b - a} \, \frac{4a + 4b}{a^2 + 2ab + b^2}
%            \item \frac{x^2 - 6x + 9}{x^2 - x - 6} \, \frac{3x + 6}{2x^2 - 7x + 3}
%            \item \frac{3x^2 + 2x - 1}{2x + 6} \div \frac{x^2 - 1}{x^2 + 2x - 3}
%            \item \frac{58}{3 (3t + 2)} + \frac{1}{3}
%            \item \frac{2x}{2x - 1} - \frac{3x}{2x + 5}
%            \item \frac{-xe^x}{x + 1} + e^x
%            \item \frac{4}{x^2 - 9} - \frac{5}{x^2 - 6x + 9}
%            \item \frac{x}{1 - x} + \frac{2x + 3}{x^2 - 1}
%            \item \frac{1 + \frac{1}{x}}{1 - \frac{1}{x}}
%            \item \frac{4x^2}{2 \sqrt{x^2 + 7}} + \sqrt{2x^2 + 7}
%            \item 6 (2x + 1)^2 \sqrt{x^2 + x} + \frac{(2x +1)^4}{2 \sqrt{x^2 + x}}
%            \item \frac{2x (x + 1)^{-1/2} - (x + 1)^{1/2}}{x^2}
%            \item \frac{(x^2 + 1)^{1/2} - 2x^2 (x^2 + 1)^{-1/2}}{1 - x^2}
%            \item \frac{(2x + 1)^{1/2} - (x + 2) (2x + 1)^{-1/2}}{2x + 1}
%        \end{mathlist}        
%    \end{exercicio}
    
    \begin{exercicio}
        Racionalize o denominador de cada expressão.
        \begin{mathlist}[2]
            \item \frac{1}{\sqrt{3} - 1}
            \item \frac{1}{\sqrt{x} - \sqrt{y}}
            \item \frac{a}{1 - \sqrt{a}}
            \item \frac{\sqrt{a} + \sqrt{b}}{\sqrt{a} - \sqrt{b}}
        \end{mathlist}    
    \end{exercicio}
    
    \begin{exercicio}
        Racionalize o numerador de cada expressão.
        \begin{mathlist}[2]
            \item \frac{\sqrt{x}}{3}
            \item \frac{\sqrt[3]{y}}{x}
            \item \frac{1 - \sqrt{3}}{3}
            \item \frac{1 + \sqrt{x + 2}}{\sqrt{x + 2}}
        \end{mathlist}        
    \end{exercicio}
    
    \begin{exercicio}
        Determine se as afirmações são verdadeiras ou falsas. Se verdadeiras, explique porquê. Se falsas, justifique através de um exemplo.
        \begin{alphalist}
            \item Se $b^2 - 4ac > 0$, então $ax^2 + bx + c = 0$, $a \neq 0$, tem duas raízes reais.
            \item Se $b^2 - 4ac < 0$, então $ax^2 + bx + c = 0$, $a \neq 0$, não tem raízes reais.
            \item \( \frac{a}{b + c} \left(\frac{a}{b} + \frac{a}{c}\right)\)
            \item \( \sqrt{(a + b) (b - a)} = \sqrt{b^2 - a^2}\) para todos os números reais $a$ e $b$.
        \end{alphalist}                                        
    \end{exercicio}
    
    \begin{exercicio}
        Encontre os valores de $x$ que satisfazem a(s) desigualdade(s).
        \begin{mathlist}
            \item - x + 3 \leq 2x + 9
            \item x - 3 > 2 \qquad \text{ou} \qquad x + 3 < -1
            \item 2x^2 \, > \, 50
        \end{mathlist}
    \end{exercicio}
    
    \begin{exercicio}
        Calcule as expressões dadas.
        \begin{mathlist}
            \item \abs{-5 + 7} + \abs{-2}
            \item \abs{2\pi - 6} - \pi
            \item \abs{\sqrt{3} - 4} + \abs{4 - 2\sqrt{3}}
        \end{mathlist}    
        \begin{answers}
            \begin{mathlist}[3]
                \item 4
                \item \pi -6
                \item 8 - 3 \sqrt{3}
            \end{mathlist}
        \end{answers}
    \end{exercicio}
    
    \begin{exercicio}
        Calcule as expressões dadas.
        \begin{mathlist}[2]
            \item \left(\frac{9}{4}\right)^{3/2}
            \item \frac{5^6}{5^4}
            \item \frac{(3 \times 2^{-3}) (4 \times 3^5)}{2 \times 9^3}
            \item \frac{3 \sqrt[3]{54}}{\sqrt[3]{18}}
        \end{mathlist}        
        \begin{answers}
            \begin{mathlist}[4]
                \item \sfrac{27}{8}
                \item 25
                \item \sfrac{1}{4}
                \item 3 \sqrt[3]{3}
            \end{mathlist}
        \end{answers}
    \end{exercicio}
    
    \begin{exercicio}
        Simplifique as expressões dadas.
        \begin{mathlist}[2]
            \item \dfrac{4(x^2 + y)^3}{x^2 + y}
            \item \dfrac{a^6 B^{-5}}{(a^3b^{-2})^{-3}}
            \item \dfrac{\sqrt[4]{16x^5yz}}{\sqrt[4]{81xyz^5}} \) \\[2mm] $x$, $y$ e $z$ positivos
            \item \left(\dfrac{3xy^2}{4x^3y}\right)^{-2} \left(\dfrac{3xy^3}{2x^2}\right)^3    
        \end{mathlist}  
        \begin{answers}
            \begin{mathlist}[4]
                \item \null\!\! 4\! \left(\! x^{2}\! +\! y \!\right)^{2}
                \item \frac{a^{15}}{B^{5} b^{6}}
                \item \frac{2x}{3z}
                \item 6 x y^{7}
              \end{mathlist}
          \end{answers}          
    \end{exercicio}
    
    \begin{exercicio}
        Fatore as expressões dadas.
        \begin{mathlist}[2]
            \item -2 \pi^2 r^3 + 100 \pi r^2
            \item 2v^3w + 22 vw^3 + 2u^2 vw
            \item 12t^3 - 6t^2 - 18t
        \end{mathlist}        
        \begin{answers}
            \begin{mathlist}[2]
                \item 2 \pi r^{2} \left(- \pi r + 50\right)
                \item 2 v w \left(u^{2} + v^{2} + 11 w^{2}\right)
                \item 6 t \left(2 t^{2} - t - 3\right)
            \end{mathlist}
        \end{answers}
    \end{exercicio}
    
    \begin{exercicio}
        Faça as operações indicadas e simplifique as expressões dadas.    
        \begin{mathlist}
            \item \frac{6x}{2(3x^2 + 2)} + \frac{1}{4(x + 2)}
            \item \frac{2}{3} \left(\frac{4x}{2x^2 - 1}\right) + 3 \left(\frac{3}{3x - 1}\right)
            \item \frac{-2x}{\sqrt{x + 1}} + 4\sqrt{x + 1}
        \end{mathlist}    
        \begin{answers}
            \begin{mathlist}[2]
                \item \frac{15 x^{2} + 24 x + 2}{4 \left(x + 2\right) \left(3 x^{2} + 2\right)}
                \item \frac{78 x^{2} - 8 x - 27}{3 \left(3 x - 1\right) \left(2 x^{2} - 1\right)}
                \item \frac{2 \left(x + 2\right)}{\sqrt{x + 1}}
            \end{mathlist}
        \end{answers}
    \end{exercicio}
    
    \begin{exercicio}
        Racionalize o denominador.
        \[\frac{\sqrt{x} - 1}{2 \sqrt{x}}\]        
    \end{exercicio}
    
\end{exercicios}

%------------------------------------------------------------------------------%

%------------------------------------------------------------------------------%
% \chapter{Revisão}
%------------------------------------------------------------------------------%

%------------------------------------------------------------------------------%
%\begin{exercicios}{Geometria}
%
%    \begin{exercicio}
%    \end{exercicio}
%    
%\end{exercicios}

%------------------------------------------------------------------------------%

%------------------------------------------------------------------------------%
% \chapter{Revisão}
%------------------------------------------------------------------------------%

%------------------------------------------------------------------------------%
\begin{exercicios}{Revisão}

    \begin{exercicio}
        Determine o custo mínimo $C$ (em reais) de um certo produto, dado que
        \[5 (C - 25) \geq 1,75 + 2,5C\]	
    \end{exercicio}
    
    \begin{exercicio}
        Determine o lucro máximo $P$ (em reais) resultante de uma certa transação, dado que
        \[6 (P - 2500) \leq 4 (P + 2400)\]	
    \end{exercicio}
        
    \begin{exercicio}
        Determine as distâncias entre os pontos
        \[\left(\frac{1}{2}, \sqrt{3}\right) \text{ e } \left(- \frac{1}{2}, 2 \sqrt{3}\right)\] 				
    \end{exercicio}
    
    \begin{exercicio}
        Determine o custo mínimo $C$ (em reais) dado que
        \[2 (1,5C + 80) \leq 2 (2,5C - 20)\]
    \end{exercicio}
    
    \begin{exercicio}
        Encontre o retorno máximo $R$ (em reais) dado que
        \[12 (2R - 320) \leq 4 (3R + 240)\]
    \end{exercicio}
    
\end{exercicios}

%------------------------------------------------------------------------------%


\writeAnswers{02-Revisao/02-Revisao-Exercicios-ans}

%------------------------------------------------------------------------------%
